%%
%% Manuscript: Microbiome-CDM
%%
%% Built against sn-jnl.cls v3.1 (Springer Nature LaTeX template, December
%% 2024). BMC journals accept this template for submission.
%% Build with:  latexmk -pdf manuscript.tex
%%
%% Target journal: Microbiome (ISSN 2049-2618), article type Methodology (switched from BMC Microbiology 2026-11-10)
%%   - Vancouver numbered citations in square brackets, e.g. [1, 3-5].
%%     The class option sn-vancouver-num loads natbib and sets
%%     \bibliographystyle{sn-vancouver-num}. Do NOT pass sort&compress to
%%     natbib and do NOT write \bibliographystyle yourself: the class does
%%     both, and a second \bibliographystyle aborts the run with
%%     "Illegal, another \bibstyle command".
%%   - Structured abstract: Background / Results / Conclusions, <= 350 words,
%%     no citations or figure references.
%%   - 3-10 keywords, supplied once via a single \keywords command.
%%   - Declarations block before the references: ethics approval, consent for
%%     publication, availability of data and materials, competing interests,
%%     funding, authors' contributions, acknowledgements.
%%

%% Microbiome requires double line spacing and line numbering in the main text (referee + lineno).
\documentclass[referee,lineno,pdflatex,sn-vancouver-num]{sn-jnl}
%% For a review/submission copy you can additionally request line numbers and
%% double spacing:  \documentclass[referee,lineno,pdflatex,sn-vancouver-num]{sn-jnl}

\usepackage{graphicx}
\usepackage{amsmath,amssymb,amsfonts,amsthm}
\input{glyphtounicode}\pdfgentounicode=1 % ligatures (fi, fl) map to Unicode, so extracted PDF text is complete
\usepackage{xcolor}
\usepackage{textcomp}
%% array: ragged-right paragraph columns in tables; the generated summary
%% table would otherwise stretch its narrow Value column.
\usepackage{array}

%% ORCID icons. sn-jnl pre-defines \orcidlogo to expect an Orcidlogo.eps file
%% that we do not ship, so undefine it before loading orcidlink; the package
%% then supplies its own scalable logo. This must happen in this order.
\makeatletter
\let\orcidlogo\@undefined
\makeatother
\usepackage{orcidlink}

%% Figures live in ./Fig
\graphicspath{{Fig/}}

%% Theorem environments (carried over from the previous template; still unused)
\theoremstyle{thmstyleone}%
\newtheorem{theorem}{Theorem}%
\newtheorem{proposition}[theorem]{Proposition}%
\theoremstyle{thmstyletwo}%
\newtheorem{example}{Example}%
\newtheorem{remark}{Remark}%
\theoremstyle{thmstylethree}%
\newtheorem{definition}{Definition}

%% Note: do NOT \usepackage{booktabs}. sn-jnl defines its own \toprule,
%% \midrule and \botrule, and booktabs would clash with them.

\begin{document}

%% ---------------------------------------------------------------------------
%% FRONT MATTER - replace every placeholder
%% ---------------------------------------------------------------------------
%% Title authored by the lead authors. The bracketed running title is
%% derived from it for page-header brevity; adjust if desired.
\title[Microbiome-CDM: a provenance-traceable data model]{Microbiome-CDM: a provenance-traceable common data model and multi-program extraction method for the human microbiome literature}

%% Author names/affiliations follow paper/publication/manuscript_bib/TITLE_PAGE.md
%% (2026-11-07 state; VERIFY-HOLDER items are still open there and must be
%% settled by the authors before submission).
%% In sn-jnl, \author* marks the corresponding author and \affil* the
%% corresponding affiliation; the e-mail follows the starred \author.
\author*[1]{\fnm{Jasmeer Singh} \sur{Kalra}\orcidlink{0000-0001-7836-6860}}\email{jkalra@uni-bonn.de}
\author[2,1]{\fnm{Martin} \sur{Hofmann-Apitius}\orcidlink{0000-0001-9012-6720}}
\author[3,2]{\fnm{Marie-Christine} \sur{Simon}\orcidlink{0000-0001-6625-1265}}

%% Authors/affiliations per author instruction (2026-11-07): affiliation 3 is
%% University of Bonn, Faculty of Agriculture, Institute of Nutrition and
%% Food Sciences (Nutritional Physiology) / BonnGut Junior Research Group
%% (Marie-Christine Simon); affiliation 2 department is the Bioinformatics
%% Department (Martin Hofmann-Apitius, Fraunhofer SCAI). Both dept VERIFYs
%% are settled. VERIFY remaining: streets/postcodes are the institutions'
%% public addresses; confirm against the official imprint pages before
%% submission.
\affil*[1]{\orgdiv{Bonn-Aachen International Center for Information Technology (b-it)}, \orgname{University of Bonn}, \orgaddress{\street{Friedrich-Hirzebruch-Allee 6}, \city{Bonn}, \postcode{53115}, \state{North Rhine-Westphalia}, \country{Germany}}}
\affil[2]{\orgdiv{Bioinformatics Department}, \orgname{Fraunhofer Institute for Algorithms and Scientific Computing (SCAI)}, \orgaddress{\street{Schloss Birlinghoven}, \city{Sankt Augustin}, \postcode{53757}, \state{North Rhine-Westphalia}, \country{Germany}}}
\affil[3]{\orgdiv{Faculty of Agriculture, Institute of Nutrition and Food Sciences, Nutrition and Microbiome}, \orgname{University of Bonn}, \orgaddress{\city{Bonn}, \postcode{53115}, \state{North Rhine-Westphalia}, \country{Germany}}}

%% BMC Microbiology requires a structured abstract with exactly these three
%% headings: Background, Results, Conclusions (<= 350 words, no citations).
%% The facts below are unchanged from the previous draft; only the paragraph
%% structure is now explicit. Number provenance (MANUSCRIPT_CLAIMS_REGISTER.csv):
%%   168,464 compendium samples = C06; 453 papers = C01; 4,006 = C07;
%%   2,505,684 mentions = M01; 334,231 candidate mappings = P01;
%%   18,272 accepted via two-model agreement = P05; 25,788 = R01;
%%   118 of 121 = X07; 90.1 = B01 (BC5CDR, PubTator3, in-domain);
%%   80.6 = B02 (BioRED, HunFlair2); 0.749 disease / 0.023 organism
%%   transfer = H04 / H05 (Spearman, BioRED F1 vs corpus grounding);
%%   77.0% vs 52.0% = N05; 56.0% supported = W03 (judge screen,
%%   DeepSeek-V4.1-Flash, CC BY/CC0 subset of 16,032 = W02).
%%  >> AUTHOR REWRITE REQUIRED: rewrite in your own voice, keep the
%%  >> facts, then delete this comment block.
\abstract{\textbf{Background:} The attributes that make microbiome sequencing
data reusable (body site, disease status, intervention, assay protocol) are
recorded inconsistently in structured submissions and survive mostly in
publication text. We present Microbiome-CDM, a method
that builds a provenance-traceable common data model (CDM) from the full text
of the 453 papers cited by a 168,464-sample gut microbiome compendium.
Separate extraction programs are kept as distinct mention sources;
mentions are normalised with two-model adjudication; findings are
catalogued with verbatim evidence; and relation triples receive a specific
label only when their sentence states the relation.

\textbf{Results:} Benchmark F1 of published recognisers predicts corpus grounding for
diseases (Spearman 0.749) but not organisms (0.023). Each of 2,505,684 mentions
and 25,788 triples is traced to its program, offsets and sentence. A blinded review found 61.8\%
(95\% CI 57.4--65.9) of sampled accepted mappings exact and 6.2\% (4.4--8.7)
wrong. After re-typing, 13,969 triples (54.2\%) carry a specific,
cue-verified relation; instead of forcing labels, the method flags 8,068
triples it cannot verify and keeps 1,210 associations
without a stated sign as unsigned. A further 4,966 specific triples derive
from catalogue findings. Against expert-curated BugSigDB, reported
directions agree in 77.0\% of determinable pairs versus a 52.0\% chance
baseline. Applied zero-shot to MiDRED's annotated microbe--disease relations,
typing reached a micro-F1 of 0.57, below fine-tuned models, but
signed answers agreed with the annotators in 88.5\%; extracted countries
agreed with GutBiomeDB's curation for the same studies in 95.8\%. Queried by genus and condition, it maps
reported directions for 22 genera in 10 conditions, which agree with
independent BugSigDB consensus in 25 of 44 comparable cells, close to chance.
On 100 sampled triples, two blinded annotators judged 70 and 79 supported or
partially supported by their sentence and agreed with each other on 87. 

\textbf{Conclusions:} A multi-program extraction method can turn the
literature behind a large microbiome compendium into a provenance-traceable
CDM. Recogniser choice has to be checked on the target corpus, because
benchmark performance transfers for diseases and not for organisms, and
abstaining where a relation cannot be verified keeps every assertion
auditable. Findings recorded only in figures and supplements are not
captured.}

%% BMC Microbiology accepts 3-10 keywords, supplied once. sn-jnl's \keywords
%% takes a single argument; the former bracket form (\keywords[...]) is not
%% supported. Abbreviations now live in the back matter.
\keywords{microbiome, common data model, metadata standards, FAIR, data
harmonisation, named entity recognition, large language models, text mining,
data provenance}

%% The previous template's "Key Points" boxed list is not part of the BMC
%% Microbiology article structure, so it has been removed. Its four claims
%% remain in the abstract and in Results, where BMC expects them.

\maketitle

%% ---------------------------------------------------------------------------
%% BODY
%% ---------------------------------------------------------------------------
%% Structure follows paper/publication/manuscript_bib/BIB_MANUSCRIPT_PLAN.md
%% (Problem Solving Protocol, 2,000-5,000 words; target ~4,750).
%% Journal policy: the text, figures, tables and references must be written by
%% the authors, not drafted by an LLM. Headings and word budgets only below.
%% Every number must trace to final_validation/MANUSCRIPT_CLAIMS_REGISTER.csv.

\section{Background}\label{sec:intro}
%% ------------------------------------------------------------------
%% CDM 3.1.0 SWITCH (Claude Code 2026-11-10; for the authors): CDM 3.1.0 is frozen
%% (= 3.0.3 + clean vocabulary, enrichment, triple repair/tiers; register CL01-CL11; Table 1 rows CL*).
%% If the lead approves describing 3.1.0 as the resource, update these rendered lines that name 3.0.3:
%%   lines 154, 336, 393, 688, 690, 734, 953, 1225, 1420 (line numbers before this note was inserted).
%% Suggested facts to state (authors' wording): built as 3.0.3, then deduplicated and validated into
%% 3.1.0: 157,815 terms after dedup, 116,755 kept [CL01]; 12,182 concepts after dedup, 11,912 kept [CL02];
%% 15,769 clean mappings [CL03]; 8,307 triples with a verified specific relation [CL05], tiers [CL06].
%% Data availability, Abstract and the Declaration of AI use must name the same version.
%% ------------------------------------------------------------------
%% ------------------------------------------------------------------
%% FRAMING NOTE 2026-11-09 (Claude Code; for the authors, prose not changed)
%% Goal: present minimal manual curation as a design choice, not a gap.
%% "No human annotation" is currently rendered 11 times (flags marked
%% FRAMING 2026-11-09 next to each). Keep it once in Methods (Stage 5
%% opening), once in Limitations and once in the AI declaration; the
%% Abstract may keep one design-property mention; cut the rest.
%% Keep the new framing defensible:
%%  1. "Orchestrator": the code is modular stages s01-s11 plus
%%     s12_orchestration, which is a fail-closed release-gate framework
%%     (gates G1-G10, read-only checks), not a one-command end-to-end runner.
%%     The build ran as separate SLURM jobs over several days, with manual
%%     remediation (migrations R1-R8, the 3.0.3 promotion); whole-build
%%     runtime is UNKNOWN (runtime_facts_2026-11-06). Describe what exists:
%%     automated stages, automated release gates, run registry.
%%  2. "Novel LLMs": the models are existing open-weight ones (Qwen3.5-4B,
%%     Gemma 4, DeepSeek-R1 / V4 / V4.1, Llama-3.1-8B, Mistral-7B); none new,
%%     none fine-tuned by the authors. The novelty is the protocol that
%%     combines and checks them.
%%  3. "Minimise", not "eliminate": the authors designed the schema, guards
%%     and pre-registered scoring, chose thresholds and fixed audit defects.
%%  4. Keep "concordance / judged validity, not accuracy" next to every judge
%%     percentage, and keep public-gold-standard F1 tied to the NER programs,
%%     not to the CDM.
%% ------------------------------------------------------------------
%% Budget: ~550 words.
%% ------------------------------------------------------------------
%% DRAFT TEXT, NOT SUBMISSION PROSE (the journal does not accept LLM-drafted
%% prose: plan §0 + the Springer Nature AI policy). Source:
%% paper/publication/manuscript_bib/INTRODUCTION_DRAFT.md (2026-11-07;
%% refined from the funnel draft in INTRODUCTION_FUNNEL_DRAFT.md;
%% facts and paragraph structure endorsed in conversation).
%% REVIEW GUIDANCE (2026-11-08, protocol-fit review): protocol articles
%% open from the bottleneck. Suggested
%% author opening claim (WRITE IN YOUR OWN WORDS -- do not ship this
%% scaffold sentence):
%%   "No standardised protocol exists for constructing a
%%   provenance-traceable common data model from the full text of a
%%   microbiome literature corpus; existing efforts deliver curated
%%   resources or knowledge bases over limited paper sets, but none
%%   carries per-row provenance back to source text at corpus scale."
%% The current paragraph 1 (scale/aggregation framing) states the
%% context; consider leading with the gap claim above, or folding it
%% into paragraph 2 where "None of these delivers..." already makes
%% the priority claim (open review flag below still applies).
%%  >> AUTHOR REWRITE REQUIRED: rewrite in your own voice, keep the
%%  >> facts and the \cite keys, then delete this comment block.
%% AUTHOR REVIEW FLAGS (2026-11-07) still open in the prose below:
%%   - "None of these delivers..." is a priority claim -- confirm it is
%%     defensible against MINERVA and BugSigDB, or soften it.
%% Citation keys map to the provisional numbers used in the drafts:
%% abdill2025=[1] goncalves2019=[2] geistlinger2023=[3]
%% langarica2025minerva=[4] chirumbolo2025minerva=[5]
%% langarica2025response=[6] hogan2024midred=[7] hong2025dimbre=[8]
%% higashi2024harmonization=[9] chen2025llmbench=[10]
%% zheng2023mtbench=[33]. Final numbering follows first appearance.
%% ------------------------------------------------------------------

The scale of human microbiome research is now set by data aggregation rather than
data generation. The most recent large-scale synthesis integrates 168,464 samples
across 482 BioProjects~\cite{abdill2025}. The sequencing reads behind all of these
studies are downloadable, but much of what makes them analysable (host species,
body site, disease status, intervention, assay protocol) is recorded
inconsistently or not at all in structured submission metadata~\cite{goncalves2019},
and survives in usable form only in the text of the underlying publications.
Cross-study comparison at this scale is therefore bounded not by data availability
but by how much prose a project can afford to read.

Manual harmonisation of that prose is the rate-limiting step. Standardising hundreds
of papers into a shared schema takes sustained expert reading, and curated
resources such as BugSigDB show what expert curation achieves, but their scope is
fixed by curation capacity rather than by
publication rate~\cite{geistlinger2023}.

Computational resources now approach the same problem from several
directions, and each contributes something the others do not. Annotated
corpora provide ground truth: MiDRED supports knowledge-base construction for
the microbiome~\cite{hogan2024midred}, and DiMB-RE annotates 14,450 entities
and 4,206 diet--microbiome relations in 165 publications~\cite{hong2025dimbre}.
Knowledge graphs provide integration and query: MINERVA maps microbe--disease
associations across the literature with a fine-tuned language model and a
visual atlas~\cite{langarica2025minerva}; MicrobiomeKG turns the supplementary
tables of microbiome papers into scored assertions available through the
Biomedical Data Translator~\cite{goetz2025microbiomekg}; KG-Microbe harmonises
microbial traits, environments and ontologies into FAIR, AI-ready
graphs~\cite{santangelo2026kgmicrobe}; and literature mining combined with
pathway data has linked gut microbiota to drug and neurotransmitter
pathways~\cite{liu2024mikg}. Sample-level resources provide reanalysis:
GutBiomeDB curates metadata for the same 168,464 samples and adds tools for
differential and meta-analysis~\cite{zhai2026gutbiomedb}. Language models have
been used to extract and harmonise sample metadata from structured
records~\cite{ikeda2025biosample,higashi2024harmonization}, although their
reliability in biomedical text processing remains under
evaluation~\cite{chen2025llmbench}.

What these resources leave open is complementary to them: the study context
(host, body site, condition, intervention, assay, sequencing platform) and the
findings stated in the prose of the papers behind a compendium, each held with
the program that extracted it, its character offsets and its verbatim sentence.
This is the information that community reporting standards ask authors to
provide, namely the STORMS checklist for human microbiome
studies~\cite{mirzayi2021storms}, the STREAMS guidelines~\cite{kelliher2025streams}
and the MIxS/MIMARKS minimum-information checklists~\cite{yilmaz2011mixs}, and
that the FAIR principles ask to be reusable~\cite{wilkinson2016fair}. For the
large body of already published studies, it can be recovered only from text.
%% REVIEW 2026-11-09 (Claude Code; each claim checked against the cited
%% abstract from Europe PMC / ACL Anthology; prose not changed):
%%  - geistlinger2023: OK. BugSigDB is manually curated, community-editable,
%%    >2,500 signatures from >600 studies.
%%  - hogan2024midred: OK. Human-annotated dataset, 3,116 microbe-disease
%%    relation annotations (it is a training dataset, not a knowledge base).
%%  - hong2025dimbre: OK. Annotated corpus (15 entity, 13 relation types;
%%    14,450 entities, 4,206 relations) plus fine-tuned models and a
%%    GPT-4o benchmark.
%%  - langarica2025minerva: OK, but note MINERVA's abstract itself claims an
%%    ontology-driven knowledge graph that "prioritizes accuracy and
%%    transparency" (fine-tuned LLM). A traceability contrast must be
%%    specific (offsets, run registry, separate program outputs), not general.
%%  - higashi2024harmonization: UNDERSTATED, and it challenges the priority
%%    claim. The preprint applies LLM metadata harmonisation to "thousands of
%%    human gut microbiome papers", >400,000 samples, reports 50% recovery of
%%    manually curated metadata and builds a unified searchable database. That
%%    is larger than this corpus (453 papers; 156,330 samples [C05]). "Has been
%%    explored" undersells it; the "None of these delivers ..." sentence must
%%    say concretely what differs (verbatim offset-anchored evidence per
%%    value, several programs kept separate, a CDM with a run registry, the
%%    computational evaluation). Also check whether the preprint has since
%%    been published (bibliography entry is bioRxiv 2024).
%%  - chen2025llmbench: supports "fine-tuned models outperform zero/few-shot
%%    LLMs in most BioNLP tasks; hallucinations and missing information".
%%    "Reliability ... remains contested" is a stretch: the paper measures
%%    limitations, it does not report a dispute.
%%  - Unsupported without a citation: "Standardising hundreds of papers ...
%%    takes sustained expert reading" (no source given).
%%  - Note for the next paragraph: the pipeline itself relies on LLMs
%%    (Qwen3.5-4B extraction, DeepSeek/Gemma/Llama/Mistral judges), so the
%%    reader will ask how this protocol handles the limitation just cited.

We built Microbiome-CDM as a pipeline rather than a single model. The pipeline
ingests the full text of each publication, converts it to section-aware structured
text, and applies an ensemble of recogniser programs (trained taggers, zero-shot
span detectors, small language models and a MeSH gazetteer) whose outputs are kept
separate instead of merged, so that every downstream row carries the identity of the
program that produced it. A guarded normalisation stage resolves mentions to
canonical terms with two-model adjudication for contested mappings, and a
schema-constrained LLM extracts entity--relation triples, each carrying a verbatim
sentence from its source paper as evidence. Because exhaustive manual annotation
of a corpus of this size is infeasible, the pipeline embeds its own evaluation:
structural checks, character-offset checks, LLM-as-judge
review~\cite{zheng2023mtbench}, a pre-specified comparison against the
expert-curated BugSigDB~\cite{geistlinger2023}, a blinded human review of a
random sample of accepted vocabulary mappings, and a blinded two-annotator check
of a random sample of relation triples.

Relative to the resources above, the contribution is the combination of four
properties in one released model. First, every mention, mapping and triple
carries the identity of the program that produced it, its character offsets and
its evidence sentence. Second, study context and findings sit in one relational
schema, so that a finding can be filtered by body site, assay or condition.
Third, the extraction abstains: a triple whose sentence does not state a
specific relation is flagged rather than given a generic one. Fourth, all 453
papers are rights-checked, so each released layer states what may be
redistributed. Table~\ref{tab:compare} sets these properties beside the
resources above.

\begin{table*}[!t]
\caption{Microbiome-CDM beside related microbiome resources. Each resource is
described by its own publication; the last column states how it and Microbiome-CDM complement each other and, where measured, how far they agree (Results).\label{tab:compare}}%
\footnotesize
\renewcommand{\arraystretch}{1.05}%
\tabcolsep=3pt%
\begin{tabular}{@{}>{\raggedright\arraybackslash}p{0.17\textwidth}>{\raggedright\arraybackslash}p{0.35\textwidth}>{\raggedright\arraybackslash}p{0.42\textwidth}@{}}
\toprule
Resource & What it provides & Complementarity with Microbiome-CDM \\
\midrule
BugSigDB~\cite{geistlinger2023} & Expert-curated differential-abundance signatures with study metadata & Reference standard for our findings (77.0\% direction agreement); we add mention-level provenance and text-only attributes \\
DiMB-RE~\cite{hong2025dimbre}, MiDRED~\cite{hogan2024midred} & Manually annotated corpora with entity and relation types; trained extraction models & Scored zero-shot: typed F1 0.25 (DiMB-RE) and 0.57 (MiDRED), below fine-tuned models; signed answers agree in 88.5\% (MiDRED). We supply corpus-scale candidates with evidence sentences \\
MINERVA~\cite{langarica2025minerva} & Literature-scale microbe--disease knowledge graph from a fine-tuned language model, with a visual atlas & Genus--condition directions agree in 72\% of comparable cells; our signed answers match 97\% of its labelled sentences. We add study-context attributes, per-row program identity and offsets \\
MicrobiomeKG~\cite{goetz2025microbiomekg} & Scored assertions from supplementary tables; available through the Biomedical Data Translator & Mines the supplements our text extraction covers least (22.1\% of supplement-sourced signatures); we cover the prose; 1 of its 42 source papers is in our corpus \\
KG-Microbe~\cite{santangelo2026kgmicrobe} & FAIR, ontology-grounded graphs of microbial traits and environments & Trait data for 86\% of our genera; its Disbiome directions agree with ours in 72\% of cells; we supply study-level context and findings \\
GutBiomeDB~\cite{zhai2026gutbiomedb} & Curated metadata and analysis tools for the 168,464 compendium samples & Same compendium, different layer: samples there, the papers here; our extracted country agrees with its curation in 95.8\% of studies, and we cover 77.1\% of its disease labels \\
PubTator3~\cite{wei2024pubtator3} & Entity and relation annotation across PubMed and PMC & One of our mention sources; we keep its output beside other programs rather than merged \\
LLM metadata extraction~\cite{ikeda2025biosample} & Language-model extraction of terms from structured sample records & We apply the same idea to full-text papers, with evidence offsets and a judge-based audit \\
\midrule
Microbiome-CDM (this work) & 453 rights-checked full-text papers; 2.5 million mentions from separate programs; mappings; catalogue; 25,788 triples, each traced to program, offsets and sentence & Strengths: provenance at row level, explicit abstention, licence discipline. Gap: small human checks (500 mappings; 100 triples, two annotators), no annotated gold standard \\
\botrule
\end{tabular}
\end{table*}

Here we describe the pipeline and the resulting resource. Build~3.0.3 of the
resource holds
162,409 canonical terms, 334,231 candidate mappings of which 314,772 pass
vocabulary guards and 18,272 are accepted by two-model agreement, and 25,788
canonical entity--relation triples, each carrying a sentence from its
source paper as evidence (for 24,157 of them located in the section text
verbatim or after normalisation). As a check on the extraction, an automated judge
screen by a single large language model (DeepSeek-V4.1-Flash, accessed through
the Fraunhofer Genie gateway) of the 16,032 triples whose cited sentence is CC BY
or CC0 licensed
classifies 56.0\% as supported, 22.6\% as partially supported, 18.2\% as not
supported and 3.3\% as conflicting; these are model judgements, reported as
concordance measures rather than accuracy claims. In the following sections, we
detail the multi-phase processing lineage, from corpus acquisition to the
common data model, and present the computational evaluation of the pipeline.
%% Number provenance -- ALL numbers in this paragraph are register rows
%% (paper/publication/final_validation/MANUSCRIPT_CLAIMS_REGISTER.csv):
%%   162,409 canonical terms = T02 (CDM 3.0.3, after review).
%%   334,231 / 314,772 / 18,272 = P01 / P02 / P05 (mapping candidates -> kept
%%     -> accepted). Do NOT quote 331,394 (unregistered mapping-evidence view,
%%     corrections-row M-01; risks mixing builds).
%%   25,788 triples = R01. "most carrying a verbatim sentence": 6.5% of
%%     triple-evidence mentions are unanchored (RESULTS_REFERENCE §6, limitations).
%%   16,032 judged = W02; 56.0/22.6/18.2/3.3% = W03 (CC BY/CC0 subset ONLY).
%%     Do NOT cite 14,363/5,780/4,618 (all-licence buckets, "internal only,
%%     do not cite" -- docs/reports/pipeline_flowchart.md).
%%   168,464 compendium samples = C06 ("168,000" is the cited paper's title
%%     figure; C05/C09 = BioProject-level linkage only, 156,330/92.8%).

\section{Methods}\label{sec:methods}

\subsection{Overview}\label{sec:overview}
%% Budget: ~350 words.
%% ---------- Figure 1 placeholder (plan §4, fig status §10: done) ----------
%% fig_overview supersedes fig1_protocol (which must NOT be used -- see plan
%% §10). Alternative icon style: fig_overview_illustrated (authors pick one).
\begin{figure}[htbp]
  \centering
  \includegraphics[width=\linewidth]{fig_overview.pdf}
  %% ------------------------------------------------------------------
  %% DRAFT CAPTION, NOT SUBMISSION TEXT (plan §0). Describes the actual
  %% content of fig_overview.pdf (make_overview_figure.py, all numbers
  %% read from the claims register at draw time):
  %%   stages 1-5; per-stage boxes: C01/C02/C07 corpus; M03/M04+M05/M06
  %%   extraction into M01 mentions; T02/T03 terms, P01/P02/P05
  %%   mappings; K01/K02 catalogue, R01 triples, Q01-node graph;
  %%   releases; evaluation band (6 checks; computed, no human
  %%   annotation). Figure generated with an AI coding assistant; see
  %%   the Declaration of AI use.
  %%  >> AUTHOR REWRITE REQUIRED: rewrite the legend in your own
  %%  >> words, then delete this comment block.
  %% ------------------------------------------------------------------
  \caption{Pipeline overview. Five stages build the common data model
  from the literature underlying a human gut microbiome compendium: (1)
  assembly and licence-checking of the cited corpus (453 BioProjects,
  447 publications, 4,006 sections); (2) multi-program entity extraction
  from kept-separate sources (trained and zero-shot recognisers, two
  small-LLM runs, a MeSH dictionary), yielding 2,505,684 mentions with
  character positions; (3) normalisation to 162,409 canonical terms and
  guarded ontology mapping (334,231 candidates, 314,772 kept, 18,272
  accepted by two-model agreement); (4) the CDM 3.0.3 relational build
  (catalogue records, evidence quotes, 25,788 relation triples); and (5)
  frozen, versioned releases with checksummed manifests. The evaluation
  band underneath is mostly computational (structural checks, evidence
  checks, model judges, curated-reference comparisons, NER benchmarks
  and research-question use cases), complemented by blinded human reviews
  of 500 sampled vocabulary mappings and 100 sampled relation triples.}
  \label{fig:overview}
\end{figure}

%% ------------------------------------------------------------------
%% DRAFT TEXT, NOT SUBMISSION PROSE (plan §0). Protocol design
%% principles requested in review (2026-11-08): 2-3 sentences
%% before Stage 1. All four principles are register-/plan-verified
%% facts about the pipeline; wording is scaffold.
%%  >> AUTHOR REWRITE REQUIRED: rewrite in your own voice, then
%%  >> delete this comment block.
%% ------------------------------------------------------------------
Four design choices run through all five stages
(Figure~\ref{fig:overview}).
\textbf{Provenance first:} every row produced downstream carries the
identity of the extraction program, the source paper, section and
character offsets, so any value can be traced back to its source
text. We record this provenance directly as columns rather than mapping it
to the W3C PROV-O ontology: the question an auditor asks here is which
program produced which row from which sentence, and these columns answer it
without a separate provenance graph. \textbf{Sources kept separate:} the four mention sources are
never merged; per-program identity is preserved into the CDM, which
allows later re-weighting without re-extraction.
\textbf{Mostly computational evaluation:} exhaustive manual annotation
of a corpus of this size is infeasible, so evaluation relies on
structural checks, model judges and concordance against curated
references, reported as concordance rather than accuracy, and is
anchored by blinded human reviews of random samples of accepted
vocabulary mappings and of relation triples.
%% FACT FLAG 2026-11-10 (Claude Code; prose not changed): as written this is not accurate. The 2026-11-07
%% triple screen sent 9,629 sentences from non-CC-BY/CC0 papers to FhGenie [W01]; the filter was applied to
%% the RESULTS. The FhGenie label screens (triple labels, 157,990 term labels) sent strings, no sentences:
%% counts in outputs/results/rights/B24_EXTERNAL_PROCESSING_ADDENDUM_2026-11-10.json. Authors: state what
%% was sent and that results are reported on CC BY/CC0 sentences only (decision on W01 disclosure: item 4).
\textbf{Licence-aware reporting:} every paper is licence-checked, and
results that depend on source sentences are reported only for papers
licensed CC~BY or CC0. One screen run outside our computing cluster (DeepSeek-V4.1-Flash on Genie,
a model service hosted by Fraunhofer, one of the authors' institutions)
received its sentences
before this filter was applied, including 9,629 sentences from papers
under other licences; its results are reported for the CC~BY/CC0 subset
only. Later external screens received either term and concept labels only
or, to resolve unsigned associations, text from CC~BY or CC0 papers only.

\subsection{Stage 1: Corpus and sectioning}\label{sec:stage1}
%% Budget: ~250 words (author draft ~200; authors may expand).
%% SOURCE: author-written Methods draft, 2026-11-08
%%   (paper/publication/manuscript_bib/methods_author_draft/
%%    METHODS_AUTHOR_DRAFT_2026-11-08.md, MESSAGE seq 96785).
%% Mechanical corrections applied on slotting (verified against
%% final_validation/MANUSCRIPT_CLAIMS_REGISTER.csv and the corpus_v2 2.2.0
%% manifest unless stated):
%%   - 453 papers = C01; 447 canonical / 6 duplicate-excluded = the 2.2.0
%%     corpus-manifest counts (Stage 30 dedup kept per authors; the earlier
%%     outline note calling this framing wrong is superseded). CDM-layer
%%     nuance (C02--C04): 447 DISTINCT publications, 5 of them shared by 11
%%     BioProjects -- state the distinction explicitly if prose ever needs it.
%%   - 4,006 sections = C07 (all SHA-256 hashes verify = C08); characters
%%     = 26,874,831 in the manifest ("26.9 million" kept as the authors'
%%     rounding).
%%   - 296 = independent_fulltext_crosschecked in the manifest (282
%%     COMPLETE / 14 PARTIAL); author's "fully cross-checkable" softened to
%%     "independently cross-checked" so the claim matches the recorded label.
%%   - 168,464 compendium samples = C06 (never "168,000"; title figure).
%%   Facts not used in this stage (kept for Results if wanted): 156,330 of
%%     168,464 compendium samples are in corpus BioProjects (92.8%) =
%%     C05/C06/C09, BioProject-level linkage only.

The foundational corpus was constructed to represent the literature
underlying large-scale human microbiome data aggregation. We seeded the
corpus acquisition using BioProject accessions~\cite{barrett2012bioproject}
(PRJDB, PRJNA, PRJEB) and digital object identifiers (DOIs) of the
literature cited in a recent synthesis of 168,464 samples~\cite{abdill2025}.
Automated ingestion of PDF and TXT formats was managed via a custom
processing script (\texttt{human\_microbiome\_corpus\_builder.py}), which
applied metadata heuristics and topic classification to yield an initial
corpus of 453 reference papers.

To ensure data integrity, a deduplication stage (stage 30) isolated 447
canonical papers and excluded 6 duplicates. Rather than treating full texts
as monolithic blocks, a section-aware parsing algorithm (stage 40) segmented
the canonical papers into 4,006 distinct sections, comprising 26.9 million
characters. Subsequent full-text validation independently cross-checked 296
of these papers against external bibliographic metadata. One paper
(HMPAPER\_000372) was retrieved only as a publisher redirect page (909
characters, no article text) and yields no catalogue records; it is kept in
the 453 for completeness, and analyses that need article text use the other
452.
%% REVIEW 2026-11-09 (Claude Code fact check; prose not changed):
%%  - "segmented the canonical papers into 4,006 sections": the CDM section
%%    table covers all 453 papers (453 distinct paper_id, 4,006 rows), not
%%    only the 447 canonical ones.
%%  - "cross-checked 296 ... against external bibliographic metadata": stage 46
%%    (fulltext_crosscheck_report.json) compared each paper's sections, title
%%    and abstract with an independent full-text source (282 COMPLETE, 14
%%    PARTIAL; 157 papers had no independent source). It is a full-text
%%    cross-check, not a metadata check.

%% NOTE (unregistered figure): "296" cross-checked papers is author Method
%% prose; it is NOT a register row. On the author-confirmation list
%% (item: "296 independently cross-checked" wording). Register it or
%% adjust the wording before submission.
\subsection{Stage 2: Entity extraction with multiple programs}\label{sec:stage2}
%% Budget: ~550 words.
%% SOURCE: author-written Methods draft (seq 96785), slotted 2026-11-08 with
%% mechanical corrections. The earlier "AUTHOR REWRITE REQUIRED" scaffold note
%% is superseded: the body below is author prose. An older duplicate fact
%% sheet block was removed with the duplicate Stage 2 heading earlier.
%%
%% CORRECTIONS APPLIED (verified against MANUSCRIPT_CLAIMS_REGISTER.csv):
%%   - Draft's "In total, Stage 2 extracted 2,505,735 mentions" REMOVED as
%%     layer-mismatched: 2,505,735 = terminology 1.0.1 table rows (M02); the
%%     CDM 3.0.3 mention count is 2,505,684 (M01). Dropped rather than
%%     silently renumbered. AUTHORS: re-add a total sentence if wanted, as
%%     "2,505,684 mentions in the final model" (M01).
%%   - 1,705,501 = frozen zenodo_v1 CSV ROW count (M03 mention view =
%%     1,705,449); branch text ties it to "raw entity rows", as shipped.
%%   - 186,024 = M05 (Ollama section-level); 366,123 = M04 (vLLM full-body);
%%     248,088 = M06 (MeSH gazetteer).
%%   - Compensating addition: offset governance (M07--M09) appended so the
%%     "precise character offsets" claim is substantiated: 2,505,628 of
%%     2,505,684 anchored; systematic shift in 75 sections detected and
%%     corrected (156,870 mentions = M08); all but 56 reproduce surface
%%     form = M09.
%%   - Model citations added mechanically (gliner/gliner2/gliner_biomed,
%%     scispacy, craft, jnlpba, bc5cdr, biobert, qwen35card, vllm, ollama,
%%     lipscomb2000mesh). patel2026microbialner cited as comparator.
%%
%% REGISTER FACTS retained here for reference:
%%   cpu_ensemble_v1     1,705,449 mentions = M03 (1,705,501 CSV rows)
%%   qwen35_vllm_body_v1    366,123 mentions = M04
%%   qwen_ollama_cell7c_v1  186,024 mentions = M05
%%   mesh_gazetteer_2026    248,088 mentions = M06
%%   Provenance spine: term_mention.extraction_run_id -> run registry
%%   extraction_run; every downstream row carries its source program.

To maximise entity recall without sacrificing provenance, we designed an
extraction pipeline that treats different named entity recognition (NER)
programs as independent annotators~\cite{patel2026microbialner}. The pipeline
deployed an ensemble of recogniser programs, whose outputs were intentionally
kept separate rather than merged, ensuring every downstream row retains the
exact identity of the program that produced it.

The extraction ensemble consisted of four distinct programmatic branches:
\begin{enumerate}
\item \textbf{CPU Section Ensemble:} a baseline configuration
(\texttt{zenodo\_v1}) comprising 7 discrete models, which identified
1,705,449 entity rows in the CDM across 2,817 body sections (the terminology
release lists 1,705,501 before loading).
\item \textbf{Small Open LLM (Section-level):} Qwen3.5-4B~\cite{qwen35card} on section-level inputs, generating 186,024 entity rows.
\item \textbf{Small Open LLM (Full-body):} Qwen3.5-4B deployed via
vLLM~\cite{qwen35card,kwon2023vllm} on full-text inputs, yielding 366,123
rows.
\item \textbf{MeSH Gazetteer:} a rule-based 2026 dictionary
search~\cite{lipscomb2000mesh} contributing 248,088 rows.
\end{enumerate}

All mentions were carried forward with precise character offsets linking
them to their verbatim source sentences. Of the 2,505,684 mentions recorded
in the final model, 2,505,628 are anchored at their offsets; a systematic
offset shift in 75 sections was detected and corrected, and all but 56
mentions reproduce their source surface form at the recorded offsets.
%% REVIEW 2026-11-09 (runtime naming, corrected): corpus stages 50-53 (findings,
%%    all 25,788 triples) and the full-body NER run (366,123, M04) used vLLM
%%    (code: pipelines/corpus_v2/scripts/stage5*.py; triple_mentions provider
%%    vllm-local). The section-level run (186,024, M05) is recorded as Ollama
%%    (src/scripts/slurm/qwen_sections_ner.sbatch; run id qwen_ollama_cell7c_v1,
%%    public in the CDM run registry). Text made runtime-neutral for that run on
%%    the authors' request (2026-11-09); do not call it vLLM unless it was re-run.
%% REVIEW 2026-11-09 (Claude Code fact check; prose not changed):
%%  - "identified 1,705,501 raw entity rows across 2,817 body sections": 2,817
%%    is the run's body-section universe; mentions occur in 2,364 sections
%%    (data audit PASS_D1 §5). Say "over a universe of 2,817 body sections" or
%%    give 2,364 as the sections with mentions.
%%  - "To maximise entity recall": recall was not measured on the corpus; a
%%    design goal is fine, but avoid implying a recall result.
%%  - 7-model split (audit-verified, not a register row): en_core_sci_lg
%%    960,719; gliner_biomed 297,115; gliner_medium 168,029; en_ner_craft_md
%%    119,092; en_ner_bc5cdr_md 80,584; en_ner_jnlpba_md 52,992;
%%    biobert_diseases 26,970.

\subsection{Stage 3: Normalisation and vocabulary mapping}\label{sec:stage3}
%% UPDATE 2026-11-10 (Claude Code; prose not changed). These layers are now FROZEN as CDM 3.1.0
%% (releases/cdm_v3/3.1.0, manifest sha256 9e5f6f10...) and REGISTERED as CL01-CL11; project-lead approval
%% pending. 3.1.0 = 3.0.3 unchanged + new tables, so the 3.0.3 numbers above remain true; cite both, never mixed.
%%  Dedup + validation (src/scripts/concept_dedup/, outputs/results/concept_dedup_2026-11-09/,
%%  clean_vocabulary_2026-11-10/):
%%   - concept ID forms normalised (Stage A, deterministic): 13,330 accepted IRIs -> 12,657 concepts
%%     (673 collapsed; all concepts 204,817 -> 191,160).
%%   - three-pass dedup (exact / fuzzy / abbreviation / SapBERT name + MiniLM context, with type,
%%     number, letter, sign, prefix, taxon and ontology guards): terms 162,409 -> 157,815 (4,594 merged;
%%     6,673 borderline pairs proposed, not applied); concepts 12,657 -> 12,182.
%%   - relevance filter (DeepSeek-V4.1-Flash via FhGenie; label + prefix + type only): terms kept
%%     116,755 (39,934 not a concept, 1,126 irrelevant; UNSURE kept); concepts kept 11,912.
%%   - accepted mappings 18,272 -> 15,769 clean rows (1,712 duplicates after merging; 544 junk terms
%%     such as author names and fragments; 347 concepts).
%%  Definitions (model-generated, Qwen3.5-4B local vLLM; term_definitions_llm_2026-11-09/,
%%  enriched_terminology_2026-11-09/): contextual definition for 120,091 terms, general 530,
%%  non-concept description 41,788 (162,394 / 162,409 covered); full forms 41,400; synonyms 87,952.
%%  Ontology's own definitions for 13,325 of 18,272 accepted mappings (BioPortal/OLS/UMLS/NCBI),
%%  4,912 model descriptions where the ontology has none. Accepted = both models agreed on ANY
%%  relation (EXACT 13,540 / BROADER 598 / NARROWER 1,313 / RELATED 2,821), not "both EXACT".
%%  Sensitivity: 74,729 terms (46.0%) are found only by en_core_sci_lg (83.2% of ENTITY-typed terms);
%%  excluding them leaves 87,680 terms and loses 0 of 16,750 accepted-mapping terms.
%% Budget: ~400 words.
%% SOURCE: author-written Methods draft (seq 96785), slotted 2026-11-08.
%% All numbers verified against the register, no corrections needed:
%%   166,012 -> 165,991 canonical terms = terminology layer L1 dedup -> L2
%%     (165,991 = T01, terminology release 1.0.1). Do NOT confuse with
%%     162,409 (T02) -- CDM 3.0.3 after review, the Introduction figure.
%%   334,231 candidates = P01; 314,772 kept = P02; 18,272 accepted =
%%     P05/P06 (16,750 terms).
%%   "Two-model adjudication" is CORRECT here (and only here): P05's
%%     Gemma-projecting + DeepSeek-R1-reviewing agreement. Never apply
%%     "two-model" to the Stage 5 judges.
%%   "2.5 million mentions (T1)" is the authors' rounding of the Stage 2
%%     total (2,505,735 terminology rows = M02 / 2,505,684 CDM mentions =
%%     M01; identical at this granularity).
%%   Taxonomy: 16 upper + 66 subcategories (terminology release 1.0.1).
%%   Vocabulary citations added mechanically: UMLS, OMOP/OHDSI, NCBI
%%     Taxonomy, BioPortal; DeepSeek-R1 and Gemma model cards.
%%   AUTHORS-CITATION-NOTE (2026-11-08): reference.bib also carries a
%%     verified Mondo entry (vasilevsky2022mondo, medRxiv 2022; venue
%%     corrected from Crossref's publisher field "openRxiv" to the
%%     posted-content institution). It is currently UNcited in the text,
%%     so it does not render under unsrt -- cite it if Mondo is named in
%%     the prose or the Zenodo metadata, or leave it out of the final
%%     citation set.
%%   OMOP version pin (2026-11-08, closes open ref item 5): the local
%%     index was built from the OMOP CDM v5 Athena bundle
%%     (outputs/results/omop_vocab/, downloaded 2026-10-23), whose
%%     VOCABULARY.csv records Mondo v2026-07-06 and HPO v2026-06-23.
%%     AUTHORS: decide whether to state these versions in the text
%%     (Methods or Data availability); they are not currently in the
%%     prose, which is acceptable if the Zenodo deposit records them.

Because raw NER outputs contain high orthographic variation, the pipeline
executed a multi-step terminology normalisation chain (T1--T5). Following
the assembly of the 2.5 million mentions (T1) and the extraction of their
contextual definitions (T3), a two-level canonicalisation process (T2)
collapsed textual variants, yielding an initial space of 166,012 canonical
terms, which was subsequently refined to 165,991 terms.

To bridge these terms to standardised ontologies (e.g.\
UMLS~\cite{bodenreider2004umls}, OMOP~\cite{hripcsak2015ohdsi}, NCBI
Taxonomy~\cite{schoch2020taxonomy}, BioPortal~\cite{whetzel2011bioportal}),
we implemented a guarded mapping strategy (T5). The pipeline generated
334,231 candidate mappings. Hard programmatic guards filtered this set
down to 314,772 biologically plausible mappings. To resolve the most
contested conceptual mappings, the pipeline used a two-model adjudication
framework (DeepSeek-R1~\cite{guo2025deepseekr1} and
Gemma~\cite{gemma2026gemma4}). Requiring both models to agree accepted
18,272 mappings, with the agreed relation (exact, broader, narrower or related)
recorded for each. Finally, a draft taxonomy
(T4) was exported, organising the terms into 16 upper categories and 66
subcategories, forming Terminology release 1.0.1.

For build 3.1.0 the vocabulary was curated further without altering the
3.0.3 tables. Concept identifiers that spell the same code differently
(for example the two MeSH IRI forms) were first collapsed
deterministically, reducing 13,330 accepted identifiers to 12,657 concepts.
Terms and concepts were then deduplicated in passes of increasing
tolerance: exact matching after normalisation, fuzzy matching restricted
to harmless single-word differences, expansion of abbreviated genus names
when exactly one genus in the corpus fits, and an embedding pass that
merges two items only if both their names (SapBERT~\cite{liu2021sapbert}) and their contextual
definitions (MiniLM sentence embeddings~\cite{liu2019sbert,minilmcard}) agree and the wording difference is explained by word
order, plural forms or an acronym. Guards block any merge across entity
types or between different numbers, single letters, signs, opposite
prefixes or accepted concepts, and two organism names are merged only when
they denote the same taxon, so that, for example, \textit{Lactobacillus}
and \textit{Lactococcus} remain distinct. Deduplication reduced the
162,409 canonical terms to 157,815 and the concepts to 12,182. A relevance
filter (DeepSeek-V4.1-Flash, given only each label and its entity type)
then removed strings that are not concepts, such as author names, figure
references and sentence fragments, or that lie outside biomedical
research, keeping 116,755 terms, 11,912 concepts and 15,769 accepted
mappings between them.

%% FACT 2026-11-10: re-typing layers [JV03 JV04 JV05 JV06 JV09]; scripts src/scripts/terms/retype_entity_terms.py,
%% src/scripts/triples/retype_triples_v2.py (v2.5), retype_merge_summary.py. New model run; lead approval pending.
Two further layers type what extraction left generic. First, each of the
53,780 kept terms that carried only the generic ENTITY type was re-typed by
a locally run language model (Qwen3.5-9B~\cite{qwen35card}, 4-bit AWQ
weights~\cite{lin2023awq}, served with vLLM on a single GPU), given the
term, its merged variants, its model definition and up to three corpus
sentences from distinct papers. A JSON schema restricts the answer to the
14 CDM types, to four proposed extension types for things the 14 cannot
hold (anatomy or body site, population or host group, measure or unit, and
geography), or to NOT\_A\_CONCEPT for words that name no scientific thing.
Second, each of the 14,867 triples that the repair pass marked as stating
no specific relation was re-typed (of the 17,481 triples without a specific,
cue-verified relation after repair, the remaining 2,614 were already flagged
as unsupported or carried a specific label without a verified cue) from up to three of its source sentences. The
model first quotes the words that relate subject and object and only then
chooses a relation, from the existing closed list extended by signed
association labels (positively or negatively associated, differentially
abundant, biomarker of, predicts, increases or decreases over time) or one
of three exits: association without direction, null finding, or not
supported. A relation counts as specific only when its quoted cue is found
verbatim in a source sentence; a null finding also requires negation
wording in the cue, and a reversed direction is overruled when the
sentence reads subject, cue, object in active voice. Both prompts were
frozen after constructed test sets on which 19 of 20 terms and 17 of 17
sentences were typed as expected; the prompt for catalogue-derived triples
(below) typed 18 of 20 constructed cases as expected.

%% FACT 2026-11-10: outputs/results/triples_final_2026-11-10/SUMMARY.json (register rows JV10-JV12 to add; lead approval pending)
%% Scripts: src/scripts/triples/resolve_unsigned_associations.py + unsigned_regate.py; claim_derived_triples.py +
%% retype_triples_v3.py --claims. External processing (DeepSeek, CC BY/CC0 text only) to be added to the B24 inventory.
Two further passes completed the relation layer. The 1,708 triples that
still stated only an unsigned association were read again with more
context: first the evidence sentence and its neighbours and, if no sign or
direction was found, the full section that holds the sentence. Text from
CC~BY or CC0 papers was sent to DeepSeek-V4.1-Flash on Fraunhofer's Genie
service (959 triples); all other triples were processed by the local
Qwen3.5-9B (749 triples), so that their text did not leave our computing
cluster. A signed label was kept only if its verbatim cue contained sign
wording other than ``risk'' and, for the section step, named at least one
of the two entities, and each answer was checked by a separate yes/no
verification. Finally, 13,046 catalogue findings whose evidence quote is
anchored in the paper and that name a taxon or feature, a comparison and a
direction were converted into candidate triples when the quote contains
the subject, the object and wording for the direction (7,424 candidates,
6,856 of them not already in the model). The local model typed each
candidate from its quote in section context, under the same checks and
verification.
%% REVIEW 2026-11-09 (Claude Code fact check; prose not changed):
%%  - "organising the terms into 16 upper categories and 66 subcategories":
%%    in terminology 1.0.1 the taxonomy is a DRAFT schema and 0 terms were
%%    classified (T4, classified_terms = 0). Category assignments exist only
%%    in CDM 3.0.3 taxonomy_annotation: 37,740 terms with an agreed category
%%    [T06] of 162,409. State which layer.
%%  - "biologically plausible mappings": the guards kept 314,772 [P02]; they
%%    check form and type, not biology. "high-confidence semantic links" /
%%    "rigorous consensus": P05 is acceptance by two-model agreement, with no
%%    person and no calibration (T5b/c NOT_CALIBRATED). Claim words to check.
%%  - Ontologies actually accepted (computed 2026-11-09, not registered):
%%    18,272 rows over 16,750 terms and 13,330 concepts; NCIt, NCBI Taxonomy,
%%    MeSH, ChEBI, SNOMED CT lead. Mondo/HPO versions: see the OMOP note above.
%%  - ACCEPTED is NOT 'both EXACT' (as the CDM data dictionary says): the code
%%    (src/scripts/vllm_termpop/deepseek_review.py) accepts any relation both
%%    models agree on. Of the 18,272 accepted rows: EXACT 13,540, BROADER 598,
%%    NARROWER 1,313, RELATED 2,821. 'accepted ... as high-confidence semantic
%%    links' therefore includes 4,732 non-equivalent concepts; say so, or cite the
%%    13,540 exact agreements separately (register the split first).

\subsection{Stage 4: Data catalogue and common data model}\label{sec:stage4}
%% Budget: ~450 words.
%% SOURCE: author-written Methods draft (seq 96785), slotted 2026-11-08 with
%% mechanical corrections.
%%
%% CORRECTIONS APPLIED:
%%   - "harmonized 40,141 semantic types": 40,141 is the row count of the
%%     CDM `type_harmonization` table (component_quality_audit/
%%     08_CDM_QUALITY.md), not a count of semantic types. Re-described as
%%     recorded type-harmonization rows. AUTHORS: confirm the wording.
%%   - "331,394 ontology mapping rows" REMOVED: corrections-row M-01 rules
%%     331,394 DO_NOT_CITE (unregistered companion view). Replaced with the
%%     register-citable figures: 162,409 canonical terms = T02 (CDM 3.0.3
%%     after review) and the P01/P02/P05 mapping counts already cited in
%%     the Introduction and Stage 3.
%%   - 44 tables / 8,490,496 rows = X01 + 08_CDM_QUALITY.md; migration
%%     additive = X05 (register CDM row: 79 change-log entries 3.0.2 -> 3.0.3; ids unique since 2026-11-10, triple-reference rows are XR01-XR06); 25,788
%%     triples = R01; graph 22,579 nodes = Q01; 2,505,684 mentions = M01.
%%   - Comparator citations added: lu2025microllm and chen2026buffers
%%     (LLM-based catalogue/extraction pipelines in adjacent domains).
%% Author draft's "(Phase A)/(Phase B) internal labels kept as shipped."

To synthesize isolated entities into contextual biological knowledge, we
executed a schema-constrained relation extraction, following comparable
LLM-based extraction pipelines in microbiology and protein
research~\cite{lu2025microllm,chen2026buffers}. A secondary LLM pass
(CATNER v2 utilising Qwen3.5-4B via vLLM) scanned the processed text to
extract complex findings, assays, cohorts, samples, and datasets. This was
followed by a deterministic cleanup (Phase A) and an LLM-assisted linking
phase (Phase B). To address missing vocabulary and unmatched relations, the
pipeline triggered gap-repair chains (G1--G10), which recorded 40,141
per-term type adjudications reconciling the entity types assigned by the
different recognisers with the canonical terminology, and linked catalogue
relations directly to the canonical terms.

The finalised architecture (Microbiome-CDM 3.0.3) was compiled into a highly
structured SQLite relational database (Figure~\ref{fig:cdm}). The migration to the final build
(3.0.3-rc1 to frozen) was strictly additive, generating a reproducible schema
of 44 tables containing 8,490,496 rows. The final model successfully
encapsulated 2,505,684 mentions, 162,409 canonical terms with their guarded
vocabulary mappings, and 25,788 canonical entity--relation triples, forming
a graph of 22,579 unique nodes.
%% REVIEW 2026-11-09 (Claude Code fact check; prose not changed):
%%  - The paragraph runs two different extractions together. The 25,788
%%    relation triples come from corpus_v2 stages 53-56 (qwen3.5:4b via
%%    Ollama; 32,713 corpus-stage canonical triples gated to 25,788 [R01]).
%%    CATNER v2 (Qwen3.5-4B, vLLM) extracted the catalogue records: 52,400
%%    records = findings 24,272, assays 8,545, sample groups 7,530, cohorts
%%    7,085, datasets 2,976, data assets 1,992 [K01, K02, K12]. The catalogue
%%    counts are not stated here at all.
%%  - G1 gap filling: 21,403 values accepted, 76,253 proposals rejected
%%    [K03, K04]; explicit non-values NOT_REPORTED 46,182 / CONFLICTING 4,586
%%    [K05, K06]; quotes 88.2% verbatim / 8.0% approximate / 3.8% not located
%%    [K08-K10] -- facts the Recommendations rely on but Stage 4 never gives.
%%  - Migration 3.0.2 -> 3.0.3 added 3 tables (41 -> 44, 39 identical, 2
%%    additive changes); 79 change-log entries [X05, CDM row].
%% ---------- Figure 2 (was plan §10 OPEN; now drawn) ----------
%% fig3_cdm_architecture.pdf now draws the 10 named modules M1-M10
%% with per-entity row counts and population status (make_bib_extra_
%% figures.py), so the former OPEN item is resolved by the current
%% script. Verify the figure still matches the released 3.0.3 build
%% before submission.
\begin{figure}[htbp]
  \centering
  %% --------------------------------------------------------------
  %% DRAFT CAPTION, NOT SUBMISSION TEXT (plan §0). Content matches
  %% fig3_cdm_architecture.pdf (make_bib_extra_figures.py): 10
  %% modules / 33 entities, row counts from the released DB;
  %% populated / partially populated / embedded / planned status
  %% per entity; link labels "anchored text" (evidence -> structure
  %% spans) and "triple terms"; every record links to PAPER and
  %% EXTRACTION_RUN (module M3).
  %% Number provenance: 44 tables / 8,490,496 rows = E-tables; no
  %% foreign keys (08_CDM_QUALITY.md) -- do not claim FK integrity.
  %% Counts of populated/partial/planned tables must be re-counted
  %% from the figure script or DB before quoting; they are not
  %% register rows.
  %%  >> AUTHOR REWRITE REQUIRED: rewrite the legend in your own
  %%  >> words, then delete this comment block.
  %% --------------------------------------------------------------
  \includegraphics[width=\linewidth]{fig3_cdm_architecture.pdf}
  \caption{Architecture of CDM 3.0.3. The model resolves 33 entity
  types across 10 modules covering source documents, document
  structure, terminology, semantic mappings, processing provenance,
  quality review, evidence, findings, research context and catalogue
  releases. Per-entity row counts and population status are shown
  (fully populated, partially populated, or embedded in another
  table). Evidence records anchor to their document-structure spans,
  and finding and triple records link to canonical terms; every
  record additionally carries its source paper and extraction run,
  which together with the checksummed releases allow any derived
  value to be traced back to source text.}
  \label{fig:cdm}
\end{figure}

\subsection{Stage 5: Computational evaluation and human spot check}\label{sec:stage5}
%% UPDATE 2026-11-10 (Claude Code; prose not changed). These layers are now FROZEN as CDM 3.1.0
%% (releases/cdm_v3/3.1.0, manifest sha256 9e5f6f10...) and REGISTERED as CL01-CL11; project-lead approval
%% pending. 3.1.0 = 3.0.3 unchanged + new tables, so the 3.0.3 numbers above remain true; cite both, never mixed.
%%  - HUMAN REVIEW (if the lead approves the exception): reviewer A annotated all 500 sampled accepted
%%    mappings (blind, pre-registered; outputs/results/benchmark/human_spotcheck_mappings_2026-11-09/):
%%    EXACT 307/497 = 61.8%, BROADER 68, NARROWER 24, RELATED 67, WRONG 31/497 = 6.2%; term is a concept
%%    485/496 = 97.8%; type correct 426/493 = 86.4%. Second reviewer (100-row overlap) and the model
%%    comparison are still pending. If reported, the section title and every "no human annotation"
%%    statement must change (H01). The 100-triple check is prepared, not done.
%%  - Catalogue judges re-run with schema-enforced verdicts (same prompt and models): invalid output
%%    Mistral 53.5% -> 0.1%, Llama 7.0% -> 0.0%; all pairs observed 0.802, kappa 0.293; PASS vs FAIL
%%    observed 0.924, kappa 0.355, PABAK 0.847 at 96% PASS; opposite verdicts 6.4% of valid pairs
%%    (judge_format_analysis_constrained_2026-11-09/). Report kappa WITH PABAK and the format result.
%%    FACT CHECK 2026-11-10: format failures do NOT explain the low kappa -- with invalid output fixed
%%    (0.1%) kappa is still 0.293 [CL11]. Gwet AC1 0.771 (3 labels) / 0.913 (PASS vs FAIL); specific
%%    agreement PASS 0.959, FAIL 0.39; 3-label PABAK 0.704 [CL15, agreement_extended.json]. Reading:
%%    PASS prevalence (kappa paradox) plus genuine disagreement on FAIL/UNCERTAIN. Do not write that
%%    kappa is low because of Mistral formatting.
%% FRAMING 2026-11-09 (Claude Code; prose not changed): Stage 5 heading: a positive form keeps the meaning (e.g. annotation-free / computational evaluation); BIB_MANUSCRIPT_PLAN.md uses this heading, update it there too.
%% Budget: ~700 words.
%% SOURCE: author-written Methods draft (seq 96785), slotted 2026-11-08 with
%% mechanical corrections.
%%
%% CORRECTIONS APPLIED (all verified against the register, RESULTS_REFERENCE
%% §3/§10/§11, the judge artifacts, and 08_CDM_QUALITY.md):
%%   - Integrity tests: "118 of 121 passing perfectly" -> "118 of 121 tests
%%     pass" (2 fail, 1 conditional) = X07/X08; what the 3 consist of = X09/X10
%%     + X07_TEST_EXCEPTIONS_2026-11-10.md (S4 BioProject-link ids; S9 missing-value
%%     codes in integer columns; N3 two unreconciled claim extractions, B12). "foreign-key links" REMOVED:
%%     the build declares NO foreign keys (08_CDM_QUALITY.md structural note;
%%     identified references were checked instead).
%%   - Dimensions 2-4 numbers verified: V01 88.7 / V02 89.2 / V03 72.1;
%%     BugSigDB overlap 47 papers = N01 (Disbiome/CTD comparisons registered
%%     but with small resolved counts, kept as the authors' framing);
%%     27 configurations = B05.
%%   - S1000: "top exact F1 score of 0.705 on the PubTator3 standard" was
%%     UNREGISTERED at freeze (RESULTS_REFERENCE §11: "Not registered:
%%     register before use"). Restructured to cite registered figures only:
%%     BC5CDR 90.1 = B01, BioRED 80.6 = B02 (PubTator3/HunFlair2), best LLM
%%     Qwen3.5-4B 3-shot 69.5/66.2 = B03/B04. AUTHORS: register S1000
%%     (PubTator3 0.705, species_metrics.json) or keep it out.
%%   - Judge rubric: the draft's "Well-formedness, Microbiome-relevance,
%%     Specificity, Entity validity (W/M/S/E)" is WRONG for the 2026-11-07
%%     evidence screen -- that rubric belongs to the separate string-level
%%     entity/triple screen (L01--L13, DeepSeek-V4-Flash). The 2026-11-07
%%     DeepSeek-V4.1-Flash evidence screen (verdicts.jsonl; REPORT.md)
%%     judged each triple against its verbatim anchor sentence on:
%%     support (the verdict), polarity_ok, predicate_justified and
%%     quote_match. Corrected accordingly.
%%   - 16,032 = W02 (CC BY/CC0 subset); 56.0/22.6/18.2/3.3% = W03.
%%     NEVER cite the all-licence buckets 14,363/5,780/4,618 (B-01..B-03,
%%     corrections CSV: internal only, do not cite).
%%   - Citations: BugSigDB (geistlinger2023), Disbiome, CTD, correlation
%%     benchmark (chen2025llmbench), MT-bench judge-methodology
%%     (zheng2023mtbench), DeepSeek-V4.1-Flash (deepseekv41card).

Because the 453-paper corpus size precludes comprehensive human
ground-truth annotation, the pipeline embeds a mostly computational
evaluation framework, complemented by two small blinded human checks. We
evaluated the pipeline along the following dimensions:

\begin{enumerate}
\item \textbf{Structural Integrity:} automated scripts verified
cryptographic hashes, text offsets and identifier linkage across the 44
SQLite tables, with 118 of 121 tests passing (2 failures, 1 conditional).

\item \textbf{Metadata Concordance:} we conducted a source-first crosscheck
using deterministic matching, successfully recovering 88.7\% of accession
numbers, 89.2\% of sequencing platforms and 72.1\% of 16S rRNA gene identifiers
natively present in the text.

\item \textbf{External Curation Comparison:} we benchmarked our automated
triple extraction against expert-curated databases, specifically testing
against BugSigDB~\cite{geistlinger2023} (using 47 overlapping papers),
Disbiome~\cite{janssens2018disbiome}, and CTD~\cite{davis2026ctd}.

\item \textbf{NER Benchmarking:} the underlying NER programs were evaluated
against public gold-standard datasets (BC5CDR and BioRED) using 27 distinct
configurations. PubTator3 achieved an exact F1 of 90.1 on BC5CDR and
HunFlair2 achieved 80.6 on BioRED, with the best small LLM (Qwen3.5-4B,
3-shot) reaching 69.5 and 66.2 respectively~\cite{wei2024pubtator3,
sanger2024hunflair2,qwen35card}.

\item \textbf{Evidence-Anchored Screen:} to assess the logical validity of
the 25,788 extracted triples, we utilised an
LLM-as-judge~\cite{zheng2023mtbench} (DeepSeek-V4.1-Flash, accessed via the
Fraunhofer Genie gateway). The model evaluated each triple against its own
verbatim anchor sentence, assessing whether the predicate is justified by
the sentence, whether the polarity is correct, and whether the quoted
evidence matches the text, yielding an overall support verdict. To adhere
strictly to licensing constraints for external API processing, we report
results only for the 16,032 triples sourced from CC BY or CC0 licensed
papers. The screen classified 56.0\% of these relations as fully supported
by the text, 22.6\% as partially supported, 18.2\% as unsupported, and
3.3\% as conflicting. We report these strictly as model concordance
judgements, not as absolute accuracy
claims~\cite{chen2025llmbench,deepseekv41card}.

%% FACT 2026-11-10: judge rubric v2.1 (src/scripts/final_validation/evidence_judge_v2.py, PROMPT_SHA256 3a3e53fc...),
%% judge_v2.sbatch, judge_v2_analysis.py, link_structural_check.py [JV01 JV02 JV07 JV08]. Lead approval pending.
\item \textbf{Catalogue Judges:} three locally run model judges
(Qwen3.5-9B, Llama-3.1-8B and Mistral-7B~\cite{qwen35card,grattafiori2024llama3,jiang2023mistral})
checked catalogue values under a decomposed rubric. Links between records
were checked structurally rather than by a model. For classification
fields each judge answered the question blind, and its answer was compared
with the stored value; for stated values a judge had to quote an evidence
span, which was accepted only if found in the source passage and covering
the value; record-level summary labels were rated as fair or unfair.
Verdicts were constrained by a JSON schema. The rubric was developed on a
1\% development split (536 records, assigned by hash of the record
identifier) and frozen before the 51,837 held-out records were judged.
Agreement is reported as Cohen's and Fleiss' $\kappa$~\cite{cohen1960kappa},
Gwet's AC1 and per-verdict specific agreement, as agreement between models
rather than accuracy.

\item \textbf{Two-annotator triple check:} two annotators with microbiome
expertise independently judged a simple random sample of 100 of the 16,032
screened triples (CC BY or CC0 sentences only), using the same SUPPORTED,
PARTIAL, NOT\_SUPPORTED and CONFLICTS definitions as the model judge and
seeing only the triple and its sentence. The sample, protocol, blank sheets
and model verdicts were fixed and checksummed before annotation, and model
verdicts were opened only after both sheets were complete. Agreement is
reported as observed agreement and Cohen's $\kappa$ for the four labels and
for supported or partially supported versus not.

\item \textbf{Gold standards and curated resources:} the relation-typing and
entity prompts were applied zero-shot, unchanged and locally (Qwen3.5-9B), to
the human-annotated DiMB-RE~\cite{hong2025dimbre} and MiDRED~\cite{hogan2024midred}
corpora and to 1,100 manually labelled microbe--disease
sentences~\cite{karkera2023microbiomere} distributed with
MINERVA~\cite{langarica2025minerva}, with gold entity pairs given and our labels
mapped to each corpus's label set. Study metadata were compared per BioProject
with GutBiomeDB's curated metadata~\cite{zhai2026gutbiomedb}, and
genus-by-condition directions with MINERVA's graph and with the Disbiome
associations~\cite{janssens2018disbiome} contained in
KG-Microbe~\cite{santangelo2026kgmicrobe}; paper and genus overlap were measured
for MicrobiomeKG~\cite{goetz2025microbiomekg} and KG-Microbe. Competitor data
were used under their own terms and processed only locally.
\end{enumerate}

\paragraph{Statistical reporting.} Every percentage is stated with its
denominator or the count it is a share of, and shares are computed over
distinct items unless stated otherwise (mention-weighted validity is labelled
as such). Proportions from sampled reviews and judge agreement carry 95\%
Wilson score intervals~\cite{wilson1927}. Agreement with curated references is
tested against a chance level, either the same comparison made across
different papers or the base rate of the compared values, by an exact
one-sided binomial test. Kept and dropped triples are compared by Fisher's
exact test with an odds ratio and a Newcombe interval for the difference, and,
because triples are nested in papers, also by a paper-cluster bootstrap;
differences in benchmark F1 are tested by a paired document bootstrap with
5,000 resamples~\cite{efron1993bootstrap}. $p$-values are Holm-adjusted within
each family of related comparisons, and results that do not survive
adjustment are described as trends. The analyses are descriptive; no
sample-size calculation applies, because the corpus is fixed by the
compendium.
%% REVIEW 2026-11-09 (Claude Code fact check; prose not changed):
%%  - ITEM 5, RIGHTS (critical): "To adhere strictly to licensing constraints
%%    for external API processing, we report results only for the 16,032..."
%%    reads as if the licence filter was applied before sending. It was not:
%%    the 2026-11-07 run sent all 25,714 sentences to the Genie gateway,
%%    9,629 of them from non-CC-BY/CC0 papers [W01]; the filter was applied
%%    afterwards, to reporting only (CHANGE_RECORD 2026-11-07; B24 addendum).
%%    The wording must match what happened; decide it with the project lead.
%%    CONTEXT (author-confirmed 2026-11-09, outputs/results/rights/B24_EXTERNAL_
%%    PROCESSING_ADDENDUM_2026-11-07.json): FhGenie runs DeepSeek-V4.1-Flash on
%%    Fraunhofer-operated hardware; access was under the first author's guest
%%    contract, which permits FhGenie use; Fraunhofer SCAI is a paper affiliation.
%%    So: an internal service of a collaborating research institution, not a
%%    commercial API. Still open: FhGenie retention policy; TDM-exception
%%    confirmation (lead/legal). Accurate wording can now say the screen ran on
%%    Fraunhofer's internally hosted LLM service under the guest affiliation; it
%%    still cannot say 'local' (Marvin/Bonn) or 'filtered before sending'.
%%  - ITEM 2: "successfully recovering 88.7% of accession numbers ...
%%    natively present in the text" frames V01-V03 as recall. They are the
%%    share of the pattern extractor's findings also in the catalogue:
%%    agreement between two methods (as Results para 2 correctly says).
%%  - ITEM 3: the 47-paper BugSigDB comparison is of catalogue findings and
%%    metadata (N01-N11), not of triples; the triple check against Disbiome,
%%    BugSigDB and CTD is a separate analysis (Q03-Q09, XR01-XR06).
%%  - The five dimensions omit two evaluation layers that Results reports:
%%    catalogue values vs evidence (literal match E01/E02; Llama-3.1-8B and
%%    Mistral-7B judges J01-J08) and the DeepSeek-V4-Flash string screen
%%    (L01-L13, 730,195 items).
%%  - ITEM 4: S1000 (PubTator3 0.705) is benchmarked but not registered.

\section{Results}\label{sec:results}
%% UPDATE 2026-11-10 (Claude Code; prose not changed). These layers are now FROZEN as CDM 3.1.0
%% (releases/cdm_v3/3.1.0, manifest sha256 9e5f6f10...) and REGISTERED as CL01-CL11; project-lead approval
%% pending. 3.1.0 = 3.0.3 unchanged + new tables, so the 3.0.3 numbers above remain true; cite both, never mixed.
%%  Triple repair (Qwen3.5-4B, closed relation list, verified cue phrase; triple_repair_2026-11-09/):
%%   of 17,717 generic triples, 3,966 now carry a specific, cue-verified relation; of 8,071 specific
%%   labels, 4,117 confirmed; overall 8,307 / 25,788 specific with a verified cue (32.2%);
%%   14,867 NO_SPECIFIC_RELATION_STATED, 1,207 NOT_SUPPORTED, 93 NULL_FINDING; 468 direction reversals.
%%   Tiers: 1,794 Tier 1 / 6,263 Tier 2 / 17,731 Tier 3. Para 6 could cite the tiers as the usable subset.
%%  Para 2: the judge kappa 0.282 should now be read with the constrained re-run (Stage 5 update).
%%   FACTS 2026-11-10 (register): kappa 0.293 after schema enforcement [CL11]; AC1 0.771 / 0.913, specific
%%   agreement PASS 0.959 vs FAIL 0.39 [CL15] -> the judges agree on PASS and disagree on FAIL.
%%   Triples: 9,621 specific after repair, 8,307 with a verified cue [CL12]; after triple dedup 23,339
%%   clean triples, 8,967 with a specific relation, i.e. 61.6% WITHOUT one [CL13, CL14; not in 3.1.0].
%%   Chemical-disease CTD support 22.4% vs chance 24.2% [XR05] = at chance; gene-disease 0 of 7 [XR06].
%%   Review item 9: present triples as a secondary, screened layer; state the 61.6% and the chance result.
%% Budget: ~900 words.
%% ------------------------------------------------------------------
%% DRAFT TEXT, NOT SUBMISSION PROSE (plan §0: the journal does not accept
%% LLM-drafted prose). Same scaffold status as the Introduction.
%% Every number has its claim ID below; DO NOT delete the IDs until
%% the rewrite is done and the IDs are re-checked against the frozen
%% register (final_validation/MANUSCRIPT_CLAIMS_REGISTER.csv).
%% Judge attribution rules (plan §12.1): name the judge model with
%% every percentage; judge checks are concordance/validity screens,
%% never accuracy; judge sanity checks exist for the Qwen judge only.
%% DELIBERATELY OMITTED, pending author sign-off (see also the open
%% author-confirmation list): the S1000 species abstract benchmark
%% (PubTator3 0.705 etc., reported in plan §11.6 but NOT yet in the
%% register -- register before use); the research-question use case
%% percentages (97.1/90.1/26--39, not register rows); the IBD use
%% case (75.0% Wilson 50.5--89.8) -- only F03 (25/44, exploratory) is
%% registered; the ILM invented-string/off-schema/output-cap rates
%% (44/53/79%, plan-attested but unregistered).
%%  >> AUTHOR REWRITE REQUIRED: rewrite in your own voice, keep the
%%  >> facts and the judge attributions, then delete this comment
%%  >> block.
%% ------------------------------------------------------------------
%% ---------- Table: pipeline summary (added 2026-11-09) ----------
%% Body GENERATED from the claims register by
%% src/scripts/figures/make_pipeline_summary_table.py (Claude Code; declare in
%% the AI statement). Every value is a register row (ids in the last column;
%% X01-X03 are the CDM rows). Re-run the script after any register change.
%%  >> AUTHOR INPUT REQUIRED: review the row labels and confirm the caption
%%  >> wording. The table is cited from the first Results paragraph
%%  >> (Table~\ref{tab:summary}).
%% CAPTION (drafted 2026-11-10 during the sn-jnl conversion; it states only
%% what the generated body contains, so no new facts are introduced).
\begin{table*}[!tbp]
\caption{Summary of the extraction pipeline and of the two frozen builds.
Quantities are given per stage and are reproducible from the deposited
release; \texttt{CL} rows are the CDM 3.1.0 layers and \texttt{JV} rows the
re-typing layers and the catalogue judges.
Percentages and judge-derived values are concordance or model-judgement
measures, not accuracy against human annotation.\label{tab:summary}}%
%% GENERATED by src/scripts/figures/make_pipeline_summary_table.py from the claims register -- do not edit;
%% re-run the script. Row labels: author review. Caption: written by the authors (in manuscript.tex).
\footnotesize
\renewcommand{\arraystretch}{0.95}%
\tabcolsep=3pt%
\begin{tabular}{@{}l>{\raggedright\arraybackslash}p{0.40\textwidth}>{\raggedright\arraybackslash}p{0.27\textwidth}l@{}}
\toprule
Stage & Quantity & Value & Register \\
\midrule
Corpus & papers / distinct publications & 453 / 447 & C01, C02 \\
 & sections (SHA-256 hashed) & 4,006 & C07, C08 \\
 & compendium samples in corpus BioProjects & 156,330 of 168,464 (92.8\%) & C05, C06, C09 \\
 & papers whose text may be redistributed & 274 of 453 & RT01, RT02 \\
\midrule
Extraction & mentions, four sources kept separate & 2,505,684 & M01 \\
 & CPU NER ensemble / Qwen3.5-4B full-body / Qwen3.5-4B section-level / MeSH & 1,705,449 / 366,123 / 186,024 / 248,088 & M03--M06 \\
\midrule
Normalisation & canonical terms (after review) & 162,409 & T02 \\
 & vocabulary mappings: candidates / kept by guards / accepted by two-model agreement & 334,231 / 314,772 / 18,272 & P01, P02, P05 \\
\midrule
Catalogue & evidence-quoted records & 52,400 & K01 \\
 & quotes found verbatim at their offsets & 88.2\% & K08 \\
\midrule
CDM 3.0.3 & tables / columns & 44 / 866 & X01, X02 (CDM) \\
 & registered extraction runs & 45 & X03 (CDM) \\
 & relation triples / with a generic relation label & 25,788 / 17,717 & R01, R02 \\
 & schema, value and provenance tests: pass / fail / conditional & 118 / 2 / 1 & X07 \\
\midrule
CDM 3.1.0 & terms after dedup / kept after relevance filter & 157,815 / 116,755 & CL01 \\
 & accepted concepts after ID normalisation / dedup / kept & 12,657 / 12,182 / 11,912 & CL02 \\
 & accepted mappings, clean rows (of which EXACT agreements in 3.0.3) & 15,769 (13,540) & CL03, CL04 \\
 & triples with a specific relation and verified cue & 8,307 & CL05 \\
 & triple tiers 1 / 2 / 3 & 1,794 / 6,263 / 17,731 & CL06 \\
 & terms with a model definition: contextual / general / descriptive & 120,091 / 530 / 41,779 & CL08 \\
\midrule
Evaluation & best exact F1, public gold standards (BC5CDR / BioRED) & 90.1 / 80.6 & B01, B02 \\
 & BugSigDB: same direction of reported findings, of n pairs (chance) & 77.0\% of 261 (52.0\%) & N05 \\
 & triples judged against a CC BY / CC0 sentence (LLM screen) & 16,032 & W02 \\
 & judged supported / partial / not supported / conflicting & 56.0\% / 22.6\% / 18.2\% / 3.3\% & W03 \\
\botrule
\end{tabular}

\begin{tablenotes}%
\item LLM-screen and BugSigDB values are concordance or model judgement, not accuracy.
\end{tablenotes}
\end{table*}

%% ---------- Figure 4 (plan §10: done) ----------
\begin{figure}[htbp]
  \centering
  %% DRAFT CAPTION. Panels verified against make_bib_figures.py:fig4:
  %% (a) DeepSeek-judged valid-and-typed share of distinct strings
  %%     per NER configuration [L01--L07];
  %% (b) catalogue-vs-BugSigDB metadata concordance with chance
  %%     baselines, fields marked with a dagger compared by SapBERT
  %%     semantic match [N07, N08];
  %% (c) share of BugSigDB signatures covered by catalogue findings
  %%     by signature source (text/table/figure/supplement) [N06];
  %% (d) BugSigDB support for taxon triples kept versus dropped by
  %%     the cleaning step [XR03; trend, does not survive Holm
  %%     correction]. Figure drawn with an AI coding assistant; see
  %%     the Declaration of AI use.
  %%  >> AUTHOR REWRITE REQUIRED for the legend wording.
  %% --------------------------------------------------------------
  \includegraphics[width=\linewidth]{fig4_results.pdf}
  \caption{Computational evaluation.
  (a)~Share of distinct entity strings judged valid and correctly
  typed, per extraction configuration (string screen, distinct
  strings). (b)~Agreement between catalogue metadata and the
  expert-curated BugSigDB over the 47 shared papers, against
  cross-paper chance baselines; daggered fields are matched by
  SapBERT semantic similarity rather than exactly.
  (c)~Catalogue coverage of BugSigDB signatures by where BugSigDB
  recorded the signature. (d)~BugSigDB support for taxon triples
  kept versus dropped by model-based cleaning. All panels report
  concordance or judged validity, not accuracy against a human gold
%% FRAMING 2026-11-09 (Claude Code; prose not changed): Figure 4 caption: KEEP 'concordance or judged validity, not accuracy'; the 'human gold standard' clause can go.
  standard.}
  \label{fig:results}
\end{figure}

%% Para 1: build + structural checks [X01 X02 X06 X07 X08 F01 F02]
The five stages produced CDM build 3.0.3, a relational model of 44
tables and 866 columns in which every extracted value carries paper,
section and character-offset provenance. Of 121 schema, value and
provenance tests, 118 pass, two fail and one is conditional, and no row
names an unregistered extraction run. Seventeen entity types are
fully populated, fifteen partially (five of them embedded in other
tables' rows), and one is planned. The build, terminology and
catalogue releases were frozen with SHA-256 manifests
(Section~\ref{sec:stage4}). Build 3.1.0 leaves the entity tables unchanged and
adds the curated layers described in Section~\ref{sec:stage3}: 116,755 kept
terms, 11,912 kept concepts and 15,769 accepted mappings between them, of
which 13,540 of the original accepted rows are exact agreements and the
remainder broader, narrower or related concepts. After re-typing, no kept
term has the generic type: 103,088 (88.3\%) carry one of the 14 CDM types,
8,856 a proposed extension type and 4,811 are flagged as not a concept.
A blinded reviewer
classified 307 of 497 randomly sampled accepted mappings (61.8\%, 95\%
Wilson interval~\cite{wilson1927} 57.4--65.9\%) as exact, 159 as broader, narrower or related
(32.0\%, 28.0--36.2\%) and 31 (6.2\%, 4.4--8.7\%) as wrong, judged 485 of
496 terms (97.8\%, 96.1--98.8\%) to be genuine concepts and 426 of 493
entity types (86.4\%, 83.1--89.2\%) to be correct; the second reviewer's
100-mapping overlap is pending. Overall, 466 of 497 mappings (93.8\%,
91.3--95.6) were not wrong. Exactness depended strongly on how much support a
term had: it rose from 39.7\% (96 of 242) for terms mentioned once or twice to
89.6\% (60 of 67) for terms mentioned more than 20 times, and from 31.9\% (45 of
141) for terms found by one extraction program to 88.4\% (38 of 43) for terms
found by six; by vocabulary it ranged from 41.7\% (NCIt) to 91.2\% (ChEBI).
Because mention counts and program counts are released with every term, users
can restrict the mappings to well-supported terms. Table~\ref{tab:summary} sets
out the five stages and the two frozen builds, with the source of every value.

%% Para 2: evidence support, deterministic + catalogue judges
%% [E01 E02 J01--J08 V01--V03; judges = Llama-3.1-8B and Mistral-7B]
Deterministic matching locates 70.5\% of catalogue text values and
75.4\% of numeric counts inside their own quote-anchored evidence
passage; the remainder are dominated by coded
values that cannot be string-checked.
%% FACT 2026-11-10: JV01 JV02 JV07 JV08; ANALYSIS.json stated_majority_fail_by_field.
Three local model judges then checked 51,837 held-out catalogue records,
and all 167,197 links between records point to existing records of the
same paper. The two most consistent judges (Llama-3.1-8B and Mistral-7B)
agree on 82.9\% of 442,772 field verdicts (95\% Wilson interval
82.8--83.0; Cohen's $\kappa$ 0.566, Gwet's AC1 0.788), and when one fails a
value the other fails it too in most cases (specific agreement 0.598). With
sampling error this small the informative quantity is the variation across
fields, not the interval. Fleiss' $\kappa$ over all three judges is
0.471, 0.535 for stated values and 0.616 for the judges' blind answers to
classification questions. The remaining disagreement sits in a few fields
where the models default to different answers, for example whether a
control-group record describes one group or a case--control design. We
therefore use the judges as screening tools rather than raters: a majority
fails 13.1\% of stated values (33,253 of 254,022; 95\% Wilson interval
12.96--13.22\%), concentrated in DNA
extraction method (68.9\%), sequencing platform (56.3\%), body site
(40.4\%) and intervention agent (40.2\%). A pattern-based extractor containing no model component finds
its accession numbers, sequencing platforms and 16S rRNA gene regions present
in our catalogue in 88.7\%, 89.2\% and 72.1\% of cases, an agreement
measure between two methods, not a recall bound.

%% Para 3: judged string validity (fig 4a) [L13 L01--L12]
%% Judge = DeepSeek-V4-Flash string screen.
%% Attribution note (2026-11-08, closes open ref item 1): GLiNER v2.5
%% models follow the GLiNER2 architecture (zaratiana2025gliner2), which
%% is why this sentence cites GLiNER2 and not GLiNER-BioMed -- the
%% GLiNER-BioMed paper (btag322) describes v1.0 checkpoints only and the
%% run's v2.5 checkpoints (Ihor/gliner-large-v2.5-biomed, HF 2026-02-18)
%% postdate it. v2.5-medium needs no checkpoint citation (community
%% release); the v2.5-biomed run reports only Disbiome-overlap counts
%% (H06) which carry no paper citation in this text.
In the large string screen (730{,}195 judged items), distinct-string
validity separates the program families clearly
(Figure~\ref{fig:results}a). The MeSH gazetteer reaches 67.3\% of distinct strings
judged valid and correctly typed, BioBERT-diseases~\cite{lee2019biobert} 60.9\% and the
trained BERT chemical tagger 59.3\%; the best zero-shot recogniser,
GLiNER v2.5 medium, reaches 43.6\%~\cite{zaratiana2024gliner,zaratiana2025gliner2}; the
small language models span 15.7\% (Llama zero-shot) to 35.5\% (Gemma
zero-shot), with Qwen3.5-4B at 34.5\% (full-body run) and 36.6\% (section-level
run), or 50.7\% when mentions are weighted instead of distinct
strings. Validity also depends on CDM type: chemicals (44.3\%), data
analysis (43.1\%) and diseases (42.8\%) screen far better than
experimental methods (14.1\%) or study design (16.1\%). For the
build-3.0.3 triples, the same screen judges 39.3\% to be well-formed
claims, with valid subject and object in 38.7\%. After relation repair and
re-typing, 13,969 of the 25,788 triples (54.2\%) carry a specific
relation whose cue phrase is found in the source sentence; 8,068 are
flagged because the pipeline could not verify them against their sentence,
1,604 carry a specific label
whose cue was not found verbatim, 1,210 state an association without
direction even when the whole section is read, 924 report a null finding
and 13 have a vague endpoint. Reading more of the paper resolved 405 of
the 1,708 unsigned associations, 85 of them only from the wider section; 1,210
stayed unsigned, and the other 93 moved to another exit or received a specific
label whose cue was not found.
The catalogue's findings add 4,966 specific, cue-verified triples from 430
papers and 848 null findings, so that 18,935 of 32,644 triples (58.0\%)
carry a specific, verified relation. We treat the remainder as abstention,
not as relations delivered: where the sentence does not state a specific
relation, the triple is flagged or kept unsigned rather than given a generic
label, so that every typed relation can be checked against its cue phrase.

In the blinded two-annotator check of 100 triples, judged as originally
extracted, the annotators agreed on
whether a triple was supported or partially supported by its sentence for 87
(87.0\%, 95\% Wilson interval 79.0--92.2; Cohen's $\kappa$ 0.66) and on all four
labels for 64 ($\kappa$ 0.46); most disagreements separated SUPPORTED from
PARTIAL for hedged sentences. Annotator A judged 35 triples supported, 35
partially supported and 30 not supported; annotator B 65, 14 and 20 (one
conflicting). The DeepSeek-V4.1-Flash screen agreed with the annotators about
as well as they agreed with each other ($\kappa$ 0.34 and 0.42 on four
labels). In an exploratory comparison with the final relation layer, 33 of 42
triples that the pipeline types as specific and cue-verified were judged
supported or partially supported by annotator A (78.6\%, 64.1--88.3) and 38
by annotator B (90.5\%, 77.9--96.2). Of the 41 triples the pipeline flagged
as unverified, annotator A judged 25 and annotator B 30 to be supported or
partially supported: the flag is conservative, and many flagged triples state
a real relation that the relation vocabulary does not cover. Verified and
flagged triples therefore differ in support less than their labels suggest
(78.6--90.5\% versus 61.0--73.2\% supported or partially supported); the cue
check confirms that a specific label is worded in the sentence, not that the
remaining triples are wrong. Expanding the closed relation vocabulary is the
most direct way to recover flagged relations in a future build. Agreed labels
from the annotators' adjudication are pending.

%% Para 4: public benchmark -> corpus transfer [B01--B05 G01--G09
%% H03 H04 H05 H06 H07]
On public gold standards~\cite{li2016bc5cdr,luo2022biored}, spanning
27 configurations of 21 systems, the best systems are PubTator3 with
exact F1 90.1 on BC5CDR~\cite{wei2024pubtator3} and HunFlair2 with
80.6 on BioRED~\cite{sanger2024hunflair2}; the best small language
model (Qwen3.5-4B, three-shot prompting) reaches 69.5 and
66.2~\cite{qwen35card}. Per BioRED type, the best disease, gene/protein
and organism systems are {\em not} trained on BioRED (HunFlair2:
86.6, 83.3 and 94.4), while the best chemical, variant and cell line
systems are trained in-domain (PubTator3: 88.8, 82.3 and 80.8).
These rankings transfer to the microbiome corpus only partially:
Spearman correlation between BioRED F1 and corpus grounding is 0.749
for diseases but 0.023 for organisms (0.2 for genes and proteins, 0.252
for chemicals, 0.112 for variants), and no configuration of the 26
corpus benchmark configurations grounds more than 41.3\% of its own
extracted strings. Public benchmark F1 therefore cannot be read as
corpus-level accuracy, and program choice has to be re-checked against
the target corpus rather than inherited from benchmark leaderboards.

%% Para 5: BugSigDB concordance (fig 4b,c) [N01--N08 F03]
Against the expert-curated BugSigDB~\cite{geistlinger2023}, over 47
shared papers and 270 signatures, 27.7\% of BugSigDB taxa are covered
by our findings (exact or same genus) against a 9.0\% cross-paper
chance level, and 23.5\% by exact NCBI taxon ID; conversely, 57.3\% of
our taxa also appear in BugSigDB signatures against 18.1\% chance.
Where a direction comparison is determinable, reported directions
agree in 77.0\% of 261 taxon pairs, against a 52.0\% base rate.
Coverage tracks where BugSigDB recorded the signature: every
text-sourced signature is covered, versus 51.4\% of table-sourced,
27.6\% of figure-sourced and 22.1\% of supplement-sourced ones
(Figure~\ref{fig:results}c). Study metadata agrees closely for 16S rRNA gene
region (97.7\%), country (94.0\%), condition (92.7\%, compared by
SapBERT), body site (84.2\%) and statistical test (85.7\%), and less
well for group size (56.2\%) and group name (54.2\%).

\paragraph{Reuse example: reported directions by genus and condition.}
As a first reuse of the model, we asked a question a microbiome researcher
would ask of it: which genera do the papers report as increased or decreased
in each condition? Querying the catalogue findings by NCBI genus and by
condition returns, for the 22 most frequently reported genera and the 10 most
frequent conditions, the number of papers reporting each direction, each
traceable to its sentence (Figure~\ref{fig:taxon}). The map recovers
well-known patterns, for example \textit{Faecalibacterium} reported as
decreased in inflammatory bowel disease (IBD) and colorectal cancer, and
\textit{Fusobacterium} as increased in colorectal cancer. To compare it with
an independent source, we took BugSigDB signatures from papers outside our
corpus and kept the 44 cells with a clear consensus (at least three
signatures, at least two thirds in one direction). Our corpus reports the
same direction in 25 of them (56.8\%; 95\% Wilson interval
42.2--70.3). This is close to chance and is not a
validation; the 19 disagreements mark genus--condition pairs on which the
literature itself is divided, and the map is a reported-findings summary,
not a meta-analysis.

\begin{figure}[htbp]
  \centering
  \includegraphics[width=0.82\linewidth]{figF_taxon_condition.pdf}
  \caption{Reported direction per genus and condition, queried from the
  catalogue findings. Colour gives the balance of papers reporting the
  genus as increased (red) or decreased (blue) for the condition; the number
  is the count of papers. A check mark or cross shows agreement or
  disagreement with the independent BugSigDB consensus (signatures from
  papers outside our corpus; at least three signatures, at least two thirds
  in one direction); 25 of 44 comparable cells agree. Reported findings,
  not a meta-analysis.}
  \label{fig:taxon}
\end{figure}

\paragraph{Gold standards and curated resources.} Applied zero-shot to the
human-annotated DiMB-RE corpus with gold entity pairs given, our relation
typing assigned the exact label and direction to 37.7\% of 2,410 gold
relations (95\% Wilson interval 35.8--39.7), and its typed extraction F1 was
0.25 (document bootstrap 0.23--0.27), below the published fine-tuned
end-to-end model (0.445); most errors were relations asserted for entity pairs
the annotators left unrelated (32\% of 10,203). Entity recognition reached an
F1 of 0.69 (0.67--0.71, overlapping spans) on DiMB-RE types with a CDM
equivalent, with 90.9\% recall for microorganisms. On MiDRED's 3,116 doubly
annotated microbe--disease relations, four-class micro-F1 was 0.57
(0.55--0.59; test split 0.55, against 0.72 for GPT-4 in context and 0.91 for a
fine-tuned model); when both the gold label and our answer were signed, they
agreed in 1,920 of 2,169 cases (88.5\%, 87.1--89.8), and our answer was ``no
relation'' for 23.2\% of true relations. On 1,100 manually labelled
microbe--disease sentences, our signed answers agreed with 449 of 463 signed
labels (97.0\%, 95.0--98.2), with abstention on 47.1\% of signed sentences.
The zero-shot steps are therefore precise when they commit to a sign but
cautious, and below fine-tuned extractors in recall and overall F1.

Against curated metadata, our text-extracted fields agreed with GutBiomeDB for
the same BioProjects on country in 205 of 214 studies (95.8\%, 92.2--97.8) and
covered 229 of 297 curated disease labels (77.1\%, 72.0--81.5); disease
matching was done by the same local model, so this is a concordance measure.
Genus-by-condition directions agreed with MINERVA's graph in 95 of 132
comparable cells (72.0\%, 63.8--78.9; chance about 50\%), with most
disagreements on probiotic genera that intervention studies report as
protective, and with the curated Disbiome associations in KG-Microbe in 53 of
74 (71.6\%, 60.5--80.6); both sources may share papers with our corpus, so
these comparisons are not independent. MicrobiomeKG shares one of its 42
source publications with our corpus and contains 26.8\% of our 321 catalogue
genera, whereas KG-Microbe holds trait data for 86.0\% of them, so these graphs
complement rather than duplicate the CDM.

%% Para 6: triples, licence-filtered judge + curated references
%% (fig 4d) [W02 W03 Q03--Q09 XR03 XR05 XR06 R03 Q01 Q02]
For relation triples, a licence-filtered evidence screen had
DeepSeek-V4.1-Flash judge 16{,}032 triples against their CC BY or CC0
source sentence: 56.0\% are supported, 22.6\% partially, 18.2\% not
supported and 3.3\% conflict with the sentence. An earlier evidence
check over the eligible CC BY/CC0 triples (16{,}132 judged) found the
source sentence to state a relationship in 86.6\% of cases but to
support the original label in only 22.0\%; after cleaning-based
relabelling, support rises from 17.2\% to 50.4\%, and kept triples are
more often supported than dropped ones (30.6\% versus 17.2\%).
Against curated references, 13.2\% of 622 resolved microbe--disease
edges appear in Disbiome~\cite{janssens2018disbiome} or BugSigDB
against 5.8\% chance, and kept edges are enriched over dropped ones
(15.5\% versus 6.4\%). BugSigDB support for taxon triples is likewise
higher for kept than dropped triples (54.3\% versus 44.3\%,
Figure~\ref{fig:results}d), a trend that does not survive Holm
correction. Chemical--disease edges resolve less convincingly (11.4\%
in CTD~\cite{davis2026ctd} versus 9.8\% chance). The resulting graph
has 22{,}579 nodes and 25{,}788 edges (cleaned: 8{,}897 and 9{,}682),
with the largest connected component covering 68.5\% (full) and 67.8\%
(cleaned) of nodes.
%% REVIEW 2026-11-09 (Claude Code fact check; all Results numbers match the
%% register; prose not changed):
%%  - Para 1 "every extracted value carries paper, section and
%%    character-offset provenance" overstates: 142 datasets have no section
%%    [K11]; 3.8% of catalogue quotes are not located [K10]; 56 mentions are
%%    unanchored [M09]; 6.5% of triple-evidence mentions are unanchored
%%    (EXHAUSTIVE_INTEGRITY_REPORT.md).
%%  - Para 6 uses XR03 (54.3 vs 44.3; renumbered 2026-11-10 from X03, the
%%    triple-reference row; X01-X10 are now CDM rows only).
%%  - Para 6 "a licence-filtered evidence screen": see the Stage 5 rights flag.
%%  - The whole section is still marked DRAFT / AUTHOR REWRITE REQUIRED above.

\section{Discussion}\label{sec:discussion}
%% Microbiome: Background, Methods, Results, Discussion (with limitations), Conclusions. The opening paragraph
%% restates results already reported above; AUTHOR REVIEW of wording (AI-assisted, 2026-11-10).
The results show that a literature corpus can be turned into a
provenance-traceable common data model whose layers can be improved
without rewriting the frozen builds: re-typing removed the generic entity
type from every kept term and raised the share of triples with a
specific, cue-verified relation to 54.2\% (58.0\% with the triples derived
from catalogue findings), while flagging the triples
it cannot verify instead of forcing a label. The
evaluation also shows where the approach stops: reliability differs by
entity type, model judges agree only moderately, and findings reported
only in figures and supplements are not captured. The recommendations
and limitations below follow from these results.

\subsection{Relation to existing resources}\label{sec:compare}
Microbiome-CDM is designed to sit alongside existing microbiome resources
rather than to replace them (Table~\ref{tab:compare}). Curated databases and
annotated corpora offer expert-level precision and ground truth;
knowledge graphs offer integration, scoring and query; sample-level databases
offer reanalysis of the sequencing data. Microbiome-CDM contributes the
study context and stated findings of the full-text literature, at the level
of the individual sentence and the individual extraction program. The pairings
are concrete. Our coverage of BugSigDB signatures is lowest for those recorded
in supplements (22.1\%), which is exactly the source MicrobiomeKG mines, so
the two are complementary. GutBiomeDB curates the 168,464 samples of the same
compendium, and our model covers the 453 papers behind them, so a sample-level
query can be joined to the sentence that reports the study design. Annotated corpora such as DiMB-RE and MiDRED served as benchmarks for our
relation typing, and our layer in turn offers them a much larger pool of
candidate sentences to annotate. Measured against these resources (Results),
the CDM's study metadata agree closely with curated values, its genus-level
directions agree with two other sources in about 72\% of comparable
cells, and its zero-shot relation typing is precise when it commits to a sign
but below fine-tuned extractors in overall F1.
Because every row names its program, any of these resources can take only the
layers it trusts, for example mappings accepted by two-model agreement or
triples with a verified cue phrase. Its weakest point, compared with curated
resources, is human-annotated ground truth: our human checks (500 mappings,
100 triples) are small samples, not an annotated corpus; we report that gap in
the Limitations.


\subsection{Recommendations}\label{sec:recommendations}
%% UPDATE 2026-11-10 (Claude Code; prose not changed). These layers are now FROZEN as CDM 3.1.0
%% (releases/cdm_v3/3.1.0, manifest sha256 9e5f6f10...) and REGISTERED as CL01-CL11; project-lead approval
%% pending. 3.1.0 = 3.0.3 unchanged + new tables, so the 3.0.3 numbers above remain true; cite both, never mixed.
%%  New evidence for recipe-style items: (1) enforce the output schema for LLM judges (format failures
%%  53.5% -> 0.1%); (2) keep program outputs separate: dropping the en_core_sci_lg-only terms removes
%%  83% of untyped terms at no mapping cost; (3) re-type generic relations with a closed list plus
%%  abstain options instead of forcing a label (only 22% of generic triples state a specific relation).
%% Budget: ~300 words.
%% REVIEW GUIDANCE (2026-11-08, protocol-fit review): Protocol papers
%% are more citable when recommendations are recipe-like -- concrete
%% configurations for named corpus sizes/types, not general rules.
%% The items below are backed by register rows already cited in the
%% Results; if the authors want prescriptive phrasing, anchor items
%% to the concrete configuration facts:
%%   - corpus profile: 453 papers / 4,006 sections [C01, C07];
%%   - program mix actually used: 7 models in the CPU ensemble +
%%     two Qwen3.5-4B runs [M03-M05] + MeSH gazetteer [M06];
%%   - license strategy: judge only the CC BY/CC0 subset
%%     (16,032 triples [W02]); say so explicitly when API-
%%     restricted [W01: 25,714 sentences screened, 9,629 excluded];
%%   - L07/L08 (Qwen 34.5/36.6% vs BERT chemical 59.3/MeSH 67.3%
%%     distinct-string validity) support "reserve LLM-judge screening
%%     for triple validation, use trained taggers for entities".
%%  >> AUTHOR REWRITE REQUIRED: rewrite in your own voice, keep the
%%  >> facts, then delete this comment block (incl. this guidance).
%% ------------------------------------------------------------------
%% DRAFT TEXT, NOT SUBMISSION PROSE (plan §0). Based on the
%% decision-guide findings (register G01--G09, H04/H05):
%%  >> AUTHOR REWRITE REQUIRED: rewrite in your own voice, keep the
%%  >> facts, then delete this comment block.
%% ------------------------------------------------------------------
For a new corpus of the same kind, we recommend:

\begin{itemize}
\item Run trained taggers AND zero-shot recognisers side by side, and
choose per entity type, not per system. On BioRED the best organism,
disease and gene recognisers were not trained on BioRED at all, so
performance on the target corpus, not on the training-domain benchmark, should decide.
\item Do not read public benchmark F1 as corpus accuracy: its
correlation with corpus grounding is 0.749 for diseases but 0.023
for organisms. Re-check chosen systems against a sample of the
target corpus before committing to a configuration.
\item Keep each program's mentions separate in the CDM instead of
merging them: when the LLM sources are removed, the majority of
accepted mappings and of terms survive, so provenance per program
supports later re-weighting without re-extraction.
%% NOTE (ablation percentages): the exact ablation figures (90.4% of
%% accepted mappings, 41.1% of terms; MeSH gazetteer adding no unique
%% terms) are in plan section 11.5/WS-4 but are NOT register rows.
%% Register them (or add a corrections-CSV row) before quoting exact
%% percentages here.
\item Prefer explicit non-values over silent imputation. NOT\_REPORTED
and CONFLICTING labels preserve ambiguity (46{,}182 and 4{,}586 values
in our catalogue) and make downstream filtering a user decision.
\item Screen candidate relation triples against their own source
sentence with an LLM judge before building a graph; cleaning raised
judge-perceived support from 17.2\% to 50.4\% and enriched curated
reference coverage.
\item Attach a character-offset quote to every extracted value from
the start; deterministic re-matching against the anchored quote is
the cheapest check available, and 88.2\% of our quotes are verbatim.
\item Licence-check every text before sending it to an external
model API and record the check; this discipline is what later allows
a licence-filtered evaluation at all.
\end{itemize}
%% REVIEW 2026-11-09 (Claude Code fact check; prose not changed):
%%  - Last item contradicts the record: the 2026-11-07 screen was run
%%    WITHOUT the licence check before sending (W01: 9,629 non-CC-BY/CC0
%%    sentences sent). As a lesson learned it is valid; as a description of
%%    what this protocol did it is not. The guidance comment above
%%    ("9,629 excluded") is wrong for the same reason.
%%    CONTEXT (author-confirmed 2026-11-09, outputs/results/rights/B24_EXTERNAL_
%%    PROCESSING_ADDENDUM_2026-11-07.json): FhGenie runs DeepSeek-V4.1-Flash on
%%    Fraunhofer-operated hardware; access was under the first author's guest
%%    contract, which permits FhGenie use; Fraunhofer SCAI is a paper affiliation.
%%    So: an internal service of a collaborating research institution, not a
%%    commercial API. Still open: FhGenie retention policy; TDM-exception
%%    confirmation (lead/legal). Accurate wording can now say the screen ran on
%%    Fraunhofer's internally hosted LLM service under the guest affiliation; it
%%    still cannot say 'local' (Marvin/Bonn) or 'filtered before sending'.
%%  - "the majority of accepted mappings and of terms survive" (ablation):
%%    90.4% / 41.1% are not register rows, and 41.1% of terms is not a
%%    majority.

\paragraph{Directions for the next build.} Four measured limitations point
to concrete extensions. First, findings reported only in tables, figures and
supplements are covered far less often than findings in the text (51.4\%,
27.6\% and 22.1\% of table-, figure- and supplement-sourced BugSigDB signatures,
against all text-sourced ones), which is why supplement-based resources such as
MicrobiomeKG complement this one; a table and figure-caption pass is the first
candidate. Second, the closed relation vocabulary misses relations worded
outside it: in the two-annotator check most flagged triples stated a real
relation, so a vocabulary expanded from the corpus's flagged triples, not
from the benchmark corpora, is the second. Third, the catalogue fields for DNA
extraction method and sequencing platform are often rejected by the judge
majority (68.9\% and 56.3\%); both draw on finite vocabularies of kits and
instruments, so a gazetteer pass combined with a field-targeted extractor is the
third. Fourth, on DiMB-RE the relation step typed many entity pairs that the
annotators left unrelated; requiring a syntactic connection between the two
entities before typing would address this, but because the problem was
observed on that benchmark, the rule would need to be fixed in advance and
evaluated on new data. None of these changes the central claim that the
resource is auditable and abstains where it cannot verify a relation.

\subsection{Limitations}\label{sec:limitations}
%% UPDATE 2026-11-10 (Claude Code; prose not changed). These layers are now FROZEN as CDM 3.1.0
%% (releases/cdm_v3/3.1.0, manifest sha256 9e5f6f10...) and REGISTERED as CL01-CL11; project-lead approval
%% pending. 3.1.0 = 3.0.3 unchanged + new tables, so the 3.0.3 numbers above remain true; cite both, never mixed.
%%  - Model-generated layers (definitions, synonyms, full forms, re-typed relations, relevance filter)
%%    are not curated; known errors (e.g. wrong synonyms such as leucine for isoleucine).
%%  - After fixing judge output format, kappa remains 0.29: low because 96% of verdicts are PASS and
%%    the two models use UNCERTAIN differently; opposite verdicts are 6.4%.
%%  - Only 51% of the original specific triple labels were confirmed by the repair pass.
%%  - Term relevance screen sent term strings (no sentences) to FhGenie (Fraunhofer-hosted, guest
%%    contract); record in the rights addendum.
%% Budget: ~300 words.
%% ------------------------------------------------------------------
%% DRAFT TEXT, NOT SUBMISSION PROSE (plan §0). Sources: the
%% disclosure comment list that was at line 531-536 (moved here),
%% component_quality_audit CRITICAL rows (CQA-C-01..M-04), and the
%% register: R02 T05 D02 F01 08_CDM_QUALITY.md (no FKs; mapping
%% restored not recomputed; 2.09% definition coverage); the
%% restrictions are register- or audit-verified facts.
%% UNREGISTERED FIGURES still in this prose (attested in
%% RESULTS_REFERENCE_2026-11-06.md sections 3/9 but NOT register
%% rows; register or cut before submission):
%%   - 98.3% of ENTITY-typed terms ground nowhere (§3);
%%   - 9 of 14 CDM types have no reference vocabulary (§9);
%%   - 2.09% definition coverage (08_CDM_QUALITY.md; D02/T02 make it
%%     3,397 of 162,409 = 2.09%, derivable but not a row).
%% The 55% (89,799 of 162,409) is derivable from register rows
%% T05/T02. 8,071 typed triples = R01 - R02.
%%  >> AUTHOR REWRITE REQUIRED: rewrite in your own voice, keep the
%%  >> facts, then delete this comment block.
%% ------------------------------------------------------------------
\begin{itemize}
\item \textbf{Limited human review.} Human review covers one blinded
reviewer's check of 500 randomly sampled accepted mappings (497 rated; second-reviewer
overlap pending) and a blinded two-annotator check of 100 randomly sampled
triples (adjudication pending). These are small samples with wide intervals,
not a human-annotated corpus; all other evaluation is computational (concordance
against chance baselines, judged validity screens and curated-reference
comparisons) and is reported as concordance, never as
accuracy~\cite{chen2025llmbench}.
\item \textbf{Relation quality.} After repair and re-typing,
13{,}969 of the 25{,}788 triples (54.2\%) carry a specific relation with
a verified cue phrase; 8{,}068 are flagged as unverified, and the
two-annotator check shows that many of these state a real relation the
relation vocabulary does not cover, so the flag marks unverified triples, not
extraction errors; 1{,}210
state an association whose direction the paper does not give; a further
1{,}604 carry a specific label whose cue phrase was not found and are not
counted as verified. Re-typed relations, directions and term types are model
output, checked on small constructed test sets (17 of 17 sentences for the
CDM triples and 18 of 20 cases for the catalogue-derived triples) but the 100-triple human check judged each triple as originally extracted,
not its re-typed label, and the 195 triples
whose direction was flagged as reversed have not been reviewed. The 4{,}966
catalogue-derived triples are built by string rules from anchored catalogue
findings and then re-typed with their sentence context; they have no
separate human check. Unsigned associations were resolved by
DeepSeek-V4.1-Flash for CC BY and CC0 papers (243 of 959 resolved) and by a
local Qwen3.5-9B model for all others (162 of 749), so resolution quality may differ between licence groups. Applied zero-shot,
our relation typing is below fine-tuned extractors on DiMB-RE (typed F1 0.25
versus 0.445) and MiDRED (micro-F1 0.57 versus 0.91).
%% FACT 2026-11-10: kappa 0.282 = frozen run [J05]; schema-enforced re-run 0.293 [CL11], AC1 0.771,
%% FAIL specific agreement 0.39 [CL15]; low kappa is not a formatting artefact.
\item \textbf{Judges, not experts.} The three catalogue judges reach
moderate agreement (Cohen's $\kappa$ 0.566 for the best pair, Fleiss'
$\kappa$ 0.471), the record-level summary labels they rate are almost
always judged fair and so carry little information, and the triple screens (two judge models, two rubrics) are not
sanity-checked to the same depth for every judge; concordance with
model judges is not concordance with experts. The judges' majority rejects
68.9\% of stated DNA extraction methods, 56.3\% of sequencing platforms,
40.4\% of body sites and 40.2\% of intervention agents, so these catalogue
fields should be treated as low confidence. Each model step was run once
at temperature 0, and run-to-run variance was not measured.
%% Model-version drift (review 2026-11-08; judge identities verified
%% against the register, L13/W02/M04/M05): version-drift scaffold below.
\item \textbf{Model-version drift.} All model judgements reported here
are tied to the exact model checkpoints used (as of 2026-11): judge
behaviour can change across model updates, and any reproduction must
record the exact model versions, as the Methods do for each screen.
%%  >> AUTHOR REWRITE REQUIRED as part of the Limitations rewrite.
\item \textbf{Documented schema gaps.} In build 3.0.3, definitions found in
the papers cover 2.09\% of terms (3{,}397 of 162{,}409, in 6{,}773 definition
records); the definitions added in build 3.1.0 are model-written; the 53{,}780 kept terms re-typed from the
generic ENTITY type carry model-assigned types and mostly ground to no
vocabulary, and four of their types are proposed extensions; 9 of
14 CDM types have no reference vocabulary; the build declares no
foreign keys (multi-valued links are modelled with explicit link
tables, so referential integrity was checked by audit, not by
declared constraints); 79 change-log entries 3.0.2 $\to$ 3.0.3 were
restored, not recomputed.
\item \textbf{Transfer and scope.} Public-benchmark to corpus
transfer was established for diseases only; for organisms, an
advantage for trained taggers on abstracts was not established
outside their training data. Evaluation is restricted to the 453
gut-microbiome papers linked to one compendium, over 4{,}006
sections; generalisation to other bodies of literature is not
established. The BugSigDB comparison shares only 47 papers with
BugSigDB, and the genus-by-condition map (25 of 44 comparable cells
agree with independent BugSigDB consensus) summarises reported findings;
it is not a meta-analysis and does not establish clinical validity. The
text of the 179 papers whose licences do not permit redistribution cannot
be deposited, so the text-processing steps can be rerun on those papers
only from licensed copies.
\end{itemize}
%% REVIEW 2026-11-09 (Claude Code fact check; prose not changed):
%%  - Keep the "No human annotation" item: it is required by the project's
%%    framing rules (H01) and by the Stage 5 design.
%%  - "79 change-log entries 3.0.2 -> 3.0.3 were restored, not recomputed":
%%    "restored, not recomputed" applies to the mapping evidence (331,394
%%    rows restored in R3), not to the change log, which records the
%%    migration steps (PROMOTE, R1-R8).
%%  - Missing limitations found 2026-11-08/09: (a) text-only extraction
%%    misses figure/supplement findings (N06: figure 27.6%, supplement 22.1%
%%    covered); (b) non-English source text: one Korean-language paper
%%    (HMPAPER_000289; 335 Korean canonical terms) and PDF-extraction
%%    artefacts (mojibake, Cyrillic look-alikes); (c) ontology concepts are
%%    deduplicated by id string only: 9,047 codes stored under two id forms,
%%    no cross-ontology merge (none duplicated among accepted mappings);
%%    (d) the 2026-11-07 rights incident (W01), if the lead decides to
%%    report it here rather than in Methods.

\section{Conclusions}\label{sec:conclusions}
%% Derived from the abstract's Conclusions and the 2026-11-10 results; AUTHOR REVIEW of wording.
A multi-program extraction pipeline can turn the literature underlying a
large microbiome compendium into a provenance-traceable common data model
in which every value carries its paper, section, character offsets and
the program that produced it. Curated layers added on top of the frozen
build give every kept term a specific type and 18,935 triples a specific,
cue-verified relation, and three local model judges screen catalogue
values with moderate agreement. Reliability remains uneven across entity
types, and all judge-based results are concordance measures rather than
accuracy against human annotation. The builds are frozen with checksummed
manifests for public deposit.

%% ---------------------------------------------------------------------------
%% BACK MATTER (Declarations)
%%   BMC Microbiology order: Ethics approval and consent to participate;
%%   Consent for publication; Availability of data and materials; Competing
%%   interests; Funding; Authors' contributions; Acknowledgements.
%%   \backmatter switches sn-jnl's unnumbered heading style; \bmhead emits
%%   the run-in declaration headings BMC prints.
%% ---------------------------------------------------------------------------
\backmatter

%% Abbreviations: BMC lists these in the back matter, alphabetical, and only
%% terms that actually occur in the text. Resource names that are not
%% abbreviations (BugSigDB, BioRED, HunFlair2, PubTator3, SapBERT, MiniLM)
%% are given as names at first mention instead and are not listed here.
\section*{Abbreviations}
%% sn-jnl's description environment sets \labelwidth to zero, which collapses
%% the gap between label and text, so the list is defined explicitly here.
\begin{list}{}{%
  \setlength{\labelwidth}{6.5em}%
  \setlength{\labelsep}{0.5em}%
  \setlength{\leftmargin}{7em}%
  \setlength{\itemsep}{0pt}%
  \setlength{\parsep}{0pt}%
  \setlength{\topsep}{2pt}%
  \renewcommand{\makelabel}[1]{\textbf{#1}\hfill}}
\item[API] application programming interface
\item[BC5CDR] BioCreative V chemical--disease relation corpus
\item[CDM] common data model
\item[CTD] Comparative Toxicogenomics Database
\item[DOI] digital object identifier
\item[F1] harmonic mean of precision and recall
\item[IBD] inflammatory bowel disease
\item[LLM] large language model
\item[MeSH] Medical Subject Headings
\item[NCBI] National Center for Biotechnology Information
\item[NER] named entity recognition
\item[OMOP] Observational Medical Outcomes Partnership
\item[UMLS] Unified Medical Language System
\end{list}

\section*{Declarations}

\bmhead{Ethics approval and consent to participate}
%% BMC requires this heading even when it does not apply.
Not applicable. This study analysed only previously published literature and
public data; no human participants, human tissue or animal subjects were
involved.

\bmhead{Consent for publication}
Not applicable.

\bmhead{Availability of data and materials}
%% UPDATE 2026-11-10 (Claude Code; prose not changed). These layers are now FROZEN as CDM 3.1.0
%% (releases/cdm_v3/3.1.0, manifest sha256 9e5f6f10...) and REGISTERED as CL01-CL11; project-lead approval
%% pending. 3.1.0 = 3.0.3 unchanged + new tables, so the 3.0.3 numbers above remain true; cite both, never mixed.
%%  Candidate addition to the deposit (after approval and a rights check of model text that may
%%  paraphrase non-redistributable papers): clean vocabulary (terms_clean, concepts_clean,
%%  term_concept_map_clean, id maps) and the triple tiers.
%% REQUIRED by BMC. Scaffold below; DOI slot remains open until the
%% Zenodo bundle (B22+B23, pending SHA-256 seal, SLURM 28715683) is
%% published. Online data and code must be cited in the reference
%% list too (add a Zenodo/software DOI reference when sealed).
%%  >> AUTHOR INPUT REQUIRED: the DOI itself, and the wording of the
%%  >> restricted-materials exception if the 179 rights-unreviewed
%%  >> papers stay withheld. Provenance: the 179 figure is attested
%%  >> in final_validation/RIGHTS_READINESS_REPORT.md (papers needing
%%  >> confirmation: 179 = 135 without licence, 29 CC BY-NC-ND, 12 CC
%%  >> BY-NC, 3 CC BY-NC-SA) but is NOT a register row; register it
%%  >> or confirm the wording before submission.
%% ---------------------------------------------------------------
%% DATA AVAILABILITY: FACT CHECKLIST (Claude Code, 2026-11-09). Facts only;
%% the authors write the statement (Springer Nature AI policy). Suggested order, one
%% short sentence each:
%%  1. WHAT is deposited: CDM build 3.0.3 (public SQLite build + column
%%     dictionary), data catalogue 3.1.1 (CSV), terminology 1.0.1 (terms,
%%     ontology mappings, definitions, taxonomy), triples with evidence
%%     offsets, corpus section hashes, per-paper licence table, NER
%%     benchmark tables, SHA-256 manifest. Source: releases/zenodo_frozen/
%%     README.md. NOT in it yet: claims register, 2026-11-03..07 evaluation
%%     results (needs a v1.1 deposit if you want to list them).
%%  2. COVERAGE: every record for all 453 papers [C01] (25,788 triples [R01],
%%     162,409 terms [T02], 52,400 catalogue records [K01]).
%%  3. TEXT HANDLING: verbatim text (sections, quotes, evidence sentences,
%%     and extracted values that reproduce the source wording) is included
%%     only for papers whose licence permits redistribution (CC BY / CC0,
%%     274 papers [RT01]); for the others the DOI and character offsets point
%%     to the source. (Optional detail: 2,103,395 values removed, rights scan
%%     0 unexplained hits [RT03]. The number 179 is not needed.)
%%  4. WHERE: Zenodo, DOI [to be added when the deposit is published].
%%     Do NOT write "are openly deposited" before the upload exists; until
%%     then use future/conditional wording or leave the DOI slot.
%%  5. LICENCE of the deposit: proposed CC BY 4.0 -- pending authors' and
%%     legal sign-off (releases/ZENODO_UPLOAD_METADATA.json).
%%  6. CODE: pipeline and analysis code repository [to be added].
%%  7. DROP: the "exception requested from the editorial office" sentence
%%     and "rights status was not reviewed" (all 453 were licence-checked).
%% ---------------------------------------------------------------
The outputs of this pipeline will be deposited on Zenodo before
publication and made available to reviewers at submission: CDM build
3.1.0 (the unchanged 3.0.3 relational database, the curated 3.1.0 layers
and checksummed manifests), the data catalogue release 3.1.1, the
terminology release 1.0.1, the frozen claims register used for number
verification, and the benchmark and evaluation artefacts, DOI [to be
added when the deposit is published]. The pipeline code will be released
in a public repository with an archived version, [repository and DOI to
be added by the authors]. The underlying source texts are third-party
publications and are redistributed only under their own licences: all
453 papers were licence-checked, and the text of the 179 papers whose
licence does not permit redistribution (135 without a licence and 44
under non-commercial or no-derivatives terms) is withheld from the
deposit and is available from the original publishers. Derived annotations
over those papers (terms, mappings, triples, character offsets and text
hashes) contain no running text, but verbatim evidence quotations do; quotations
are included only for the 274 CC BY or CC0 papers, and the release of derived
annotations for the 135 papers without a licence is under legal review.
%% REVIEW 2026-11-09 (Claude Code fact check; prose not changed):
%%  - "are openly deposited": nothing is uploaded yet (releases/zenodo_frozen
%%    is built, rights scan PASS, NOT uploaded; licence and legal sign-off
%%    pending). Project rule: no public-availability statement until it
%%    exists.
%%  - "the 179 papers whose rights status was not reviewed": all 453 were
%%    reviewed (PAPER_RIGHTS_MATRIX.csv): 274 INCLUDE_TEXT, 15
%%    RESTRICTED_NC_LAYER, 164 WITHHOLD_TEXT. 179 = 164 + 15 papers whose
%%    licence does not allow redistribution (or none was found).
%%  - zenodo_frozen holds the CDM (public, text-redacted), catalogue,
%%    terminology, triples, corpus hashes, rights and validation/NER tables;
%%    it does NOT yet contain the claims register or the 2026-11-03..07
%%    evaluation results (a v1.1 deposit would be needed).


\bmhead{Competing interests}
%% >> AUTHOR INPUT REQUIRED (the visible marker below is deliberate: it must
%% >> survive into the PDF until an author replaces it). Conflicts are author
%% >> facts and cannot be inferred, so the block is not pre-filled.
%% Option A (no conflicts -- use only if all three authors confirm):
%% None declared.
%% Option B (disclosure -- keep only the lines that apply):
%% \textbf{J.S.K.} declares [conflict]. \textbf{M.H.-A.} declares
%% [conflict]. \textbf{M.-C.S.} declares [conflict].
%% If a funder had any role in study design, data collection, analysis,
%% interpretation or the decision to submit, state that role here.
\textbf{[AUTHOR INPUT REQUIRED --- per-author competing-interests statement
missing.} Replace this paragraph with ``The authors declare that they have
no competing interests.'' if that is true for all three authors, or with the
individual disclosures.\textbf{]}


\bmhead{Funding}
%% CANNOT SHIP AS HIDDEN PLACEHOLDER: grant numbers are author facts the
%% assistant cannot know or invent, so the marker below stays visible in the
%% PDF until an author replaces it. The BMC declaration order places Funding
%% after Competing interests and before Authors' contributions.
%% >> AUTHOR INPUT REQUIRED: replace the bracketed paragraph with one of
%% >> "This work was supported by <funder> [grant number <xxx>]."
%% >> "This research received no specific grant from any funding agency in
%% >>  the public, commercial or not-for-profit sectors."
\textbf{[AUTHOR INPUT REQUIRED --- funding statement missing.} Replace this
paragraph with the funder and grant numbers, or with ``This research received
no specific grant from any funding agency in the public, commercial or
not-for-profit sectors.''\textbf{]}

\bmhead{Authors' contributions}
%% >> AUTHOR INPUT REQUIRED: CRediT roles per author (BMC expects the
%% >> CRediT taxonomy, https://credit.niso.org/). Role assignment is an
%% >> author fact and is not pre-filled; the bracketed marker below stays
%% >> visible in the PDF until an author replaces it. All authors must
%% >> approve the finished statement.
%% CRediT roles: Conceptualization; Data curation; Formal analysis;
%% Funding acquisition; Investigation; Methodology; Project
%% administration; Resources; Software; Supervision; Validation;
%% Visualization; Writing -- original draft; Writing -- review \& editing.
\textbf{[AUTHOR INPUT REQUIRED --- CRediT statement missing.} Replace this
paragraph with one CRediT role list per author, drawn from the taxonomy at
\nolinkurl{https://credit.niso.org/}.\textbf{]}
%% Example of a filled line (do not ship this, it is illustrative):
%% J.S.K.: Conceptualization, Methodology, Software, Investigation,
%%        Formal analysis, Data curation, Visualization,
%%        Writing -- original draft.
%% All authors read and approved the final manuscript.

\bmhead{Acknowledgements}
%% BMC requires an "Acknowledgements" heading (British spelling); in the BMC
%% declaration order it follows Authors' contributions and precedes the AI-use
%% declaration. Heading spelling fixed 2026-11-08.
%% >> AUTHOR INPUT REQUIRED: names of people and facilities to thank,
%% >> including compute or storage allocations. Acknowledgements are
%% >> author facts; the bracketed marker below stays visible in the PDF
%% >> until an author replaces it.

\textbf{[AUTHOR INPUT REQUIRED --- acknowledgements missing.} Replace this
paragraph with the individuals, facilities and compute or storage allocations
to be acknowledged, or with ``Not applicable.'' if there are none.\textbf{]}

\bmhead{Declaration of generative AI and AI-assisted technologies in the
writing and research process}
%% UPDATE 2026-11-10 (Claude Code; prose not changed). These layers are now FROZEN as CDM 3.1.0
%% (releases/cdm_v3/3.1.0, manifest sha256 9e5f6f10...) and REGISTERED as CL01-CL11; project-lead approval
%% pending. 3.1.0 = 3.0.3 unchanged + new tables, so the 3.0.3 numbers above remain true; cite both, never mixed.
%%  Add if these layers are reported: Qwen3.5-4B for definitions / synonyms / full forms / relation
%%  re-typing (local vLLM); DeepSeek-V4.1-Flash (FhGenie) for the term and concept relevance screens;
%%  SapBERT and MiniLM embeddings for dedup; constrained re-run of Llama-3.1-8B / Mistral-7B judges.
%% REQUIRED by BMC/Springer Nature: declare here (or in Methods) AND in the
%% cover letter AND in the supplementary material. The block below is a
%% DRAFT SCAFFOLD assembled from AI_ASSIST_LOG.md (complete through
%% 2026-11-08) and register rows; authors verify every line.
%%  >> AUTHOR REWRITE REQUIRED: confirm wording, mirror it in the cover
%%  >> letter and supplementary declaration, then delete this comment.
%%
%% AUTHORS: language and style checking of this manuscript was
%% performed with an AI assistant (opencode, GLM-5.3 model, 2026-11-08:
%% banned-vocabulary, template-pattern, sentence-rhythm and number-
%% register checks; no text was rewritten by the assistant). Declare
%% these language checks in this Declaration or in Acknowledgements,
%% as Springer Nature permits for declared language checking.
%%
%% Judge-model attribution (register-verified; do not mix up):
%%   - string screen L01--L13: DeepSeek-V4-Flash, 730,195 items [L13];
%%   - 2026-11-07 triple evidence screen: DeepSeek-V4.1-Flash via the
%%     Fraunhofer Genie gateway, 16,032 triples vs CC BY/CC0 sentences
%%     [W02]; all 25,714 sentences were SENT and judged, 9,629 of them
%%     from other licences, excluded from REPORTING only [W01] (corrected 2026-11-10);
%%   - mapping adjudication (P05, 18,272 accepted rows): two models,
%%     Gemma and DeepSeek-R1;
%%   - catalogue value judges: Llama-3.1-8B [J01] and Mistral-7B [J02];
%%   - judge sanity checks: Qwen3.5-9B only (percentages not register
%%     rows; keep qualitative or register first);
%%   - extraction components: Qwen3.5-4B [M04 vLLM body 366,123;
%%     M05 Ollama 186,024]; benchmark LLMs: Gemma, Llama, Qwen 3-shot.
During the preparation of this work the
authors used generative AI in two distinct roles.
\textbf{(1) As pipeline components and evaluation instruments} (Methods,
Stages 2--5): the small language models Qwen3.5-4B (two extraction runs),
Gemma, Llama and Qwen with few-shot prompting served as entity- and
relation-extraction systems under evaluation; model judges were used as
computational screens, namely DeepSeek-V4-Flash for the string-validity
screen (730{,}195 judged items), DeepSeek-V4.1-Flash, accessed through the
Fraunhofer Genie gateway, for the licence-filtered triple evidence screen
(16{,}032 triples judged against their Creative Commons CC BY or CC0 source
sentences, drawn from 25{,}714 licence-screened sentences), Gemma and
DeepSeek-R1 as the two adjudicating models for vocabulary mapping, and
Qwen3.5-9B, Llama-3.1-8B and Mistral-7B, run locally, as judges of
catalogue values. Where the
manuscript reports a percentage from these screens, the responsible judge
model is named alongside it. All judge-derived results are reported as
concordance or judged validity, not as accuracy against human annotation;
human evaluation consists of the blinded review of 500 sampled accepted
mappings and the blinded two-annotator check of 100 sampled triples. Qwen3.5-9B, run locally, re-typed generic terms and unspecific
triples and typed triples derived from catalogue findings (Methods); to
resolve unsigned associations, DeepSeek-V4.1-Flash on the Genie service
read text from CC~BY or CC0 papers only and the local Qwen3.5-9B read all
other text. For build 3.1.0, Qwen3.5-4B also generated contextual
definitions and repaired triple relations, SapBERT and MiniLM embeddings
supported deduplication, and DeepSeek-V4.1-Flash screened term and concept
labels for relevance.
\textbf{(2) In manuscript and code preparation:} an AI coding assistant
(Claude Code, and opencode with the GLM-5.3 model) was used to develop
pipeline and analysis code, to verify every number in this manuscript
against a frozen claims register, to convert the authors' Methods draft to
this \LaTeX\ template mechanically, and to produce internal drafting
scaffolds that the authors re-wrote in their own words before inclusion;
figures were produced with the same assistance. The authors reviewed and
edited all AI-assisted outputs and take full responsibility for the content
of this published article.
%% REVIEW 2026-11-09 (Claude Code fact check; prose not changed):
%%  - "drawn from 25,714 licence-screened sentences" and the comment above
%%    ("9,629 ... excluded before scoring [W01]") are not accurate: all
%%    25,714 sentences were sent and judged; the 9,629 non-CC-BY/CC0 ones
%%    were excluded from reporting only. See the Stage 5 rights flag.
%%    CONTEXT (author-confirmed 2026-11-09, outputs/results/rights/B24_EXTERNAL_
%%    PROCESSING_ADDENDUM_2026-11-07.json): FhGenie runs DeepSeek-V4.1-Flash on
%%    Fraunhofer-operated hardware; access was under the first author's guest
%%    contract, which permits FhGenie use; Fraunhofer SCAI is a paper affiliation.
%%    So: an internal service of a collaborating research institution, not a
%%    commercial API. Still open: FhGenie retention policy; TDM-exception
%%    confirmation (lead/legal). Accurate wording can now say the screen ran on
%%    Fraunhofer's internally hosted LLM service under the guest affiliation; it
%%    still cannot say 'local' (Marvin/Bonn) or 'filtered before sending'.
%%  - "internal drafting scaffolds that the authors re-wrote in their own
%%    words before inclusion": true only once the sections marked AUTHOR
%%    REWRITE REQUIRED (Introduction, Results, Recommendations, Limitations,
%%    figure legends, this Declaration) have actually been rewritten. The
%%    declaration must match the submitted text.
%%  - Also used, to declare: DeepSeek-V4.1-Flash (Genie) label-translation
%%    check of 117 vocabulary labels (2026-11-08); Claude Code skills for
%%    style and overlap checks (AI_ASSIST_LOG.md).


%% ---------------------------------------------------------------------------
%% REFERENCES
%%   BMC Microbiology uses Vancouver numbered citations: numbers in square
%%   brackets, cited in order of first appearance. The class option
%%   sn-vancouver-num already loads natbib in numerical mode and sets
%%   \bibliographystyle{sn-vancouver-num}, so no \bibliographystyle command
%%   appears below -- a second one aborts the build.
%%   Entries are GENERATED from the registry-verified bibliography
%%   (submission_readiness/references_source/bibliography.json, re-verified
%%   2026-11-07) by make_reference_bib.py -- do not hand-edit reference.bib.
%%   Cite with \cite{key}; keys are the bibliography.json keys
%%   (e.g. \cite{abdill2025}). The reference-number order will follow the
%%   order of first appearance once the section text is in.
%% ---------------------------------------------------------------------------
%% CITATION GAPS (Claude Code, 2026-11-09; authors decide and place \cite):
%% reference.bib has 68 registry-verified entries; uncited ones do not
%% render. (A) USED in the pipeline but NOT cited yet:
%%   Llama/Mistral catalogue judges -> grattafiori2024llama3, jiang2023mistral
%%   scispaCy models (Stage 2) -> neumann2019scispacy (+ cohen2017craft,
%%     kim2004jnlpba, pyysalo2015bionlp13 as their training corpora)
%%   BioBERT / SciBERT taggers -> lee2019biobert, beltagy2019scibert
%%   GLiNER (v1) -> zaratiana2024gliner; HunFlair2's framework -> akbik2019flair
%%   Phi-3.5 and AWQ (benchmark LLMs) -> abdin2024phi3, lin2023awq
%%   DeepSeek-V4-Flash string screen -> deepseek2026v4, deepseekv4flashcard
%%   bootstrap CIs -> efron1993bootstrap; Wilson intervals -> wilson1927
%%   SapBERT matching (Fig. 4b caption) -> liu2021sapbert   [new]
%%   G6/G7 embeddings -> liu2019sbert, minilmcard
%%   OLS / NCBI Gene / HPO / Mondo -> cote2006ols, brown2015gene,
%%     gargano2024hpo, vasilevsky2022mondo; SRA -> leinonen2011;
%%   MIxS (if catalogue fields are aligned) -> yilmaz2011mixs
%% (B) NEW CANDIDATES found 2026-11-09 (web search, then Crossref/DataCite):
%%   luoma2023s1000   S1000 corpus -- only if the S1000 result is reported
%%   wilkinson2016fair FAIR principles -- Introduction / Discussion framing
%%   qi2023gutmdisorder gutMDisorder v2.0 -- curated databases (Introduction)
%%   karkera2023microbiomere language models for microbiome-disease
%%                    relation mining (BMC Bioinformatics) -- related work
%%   microbiorel2025  MicrobioRel generative RE, arXiv preprint -- related work
%%   medjudge2026     LLM-as-a-judge in healthcare, scoping review (arXiv
%%                    preprint) -- Stage 5 / Limitations (judge reliability)
%%   Already in the list and uncited: li2025lore, collins2026biotriplex,
%%     yazdani2026glinerbiomed (GLiNER-BioMed; checkpoint question open).
\bibliography{reference}

%% Author biographies: not part of the BMC article structure; kept commented
%% out. BMC supplies ORCID iDs in the front matter instead.
%\begin{biography}{}{\author{Jasmeer Singh Kalra.} ...}\end{biography}

%% Appendices (optional)
%\begin{appendices}
%\section{Additional derivations}\label{app:a}
%Appendix material goes here.
%\end{appendices}

\end{document}